\section{Proofs}
\label{app:proofs}


\subsection{Proof generalized conformal risk control with exchangeable data.}
\label{app:gcrc_proof_nonmonotonic_1dim}

\textbf{Notation:}
\begin{itemize}
    \item \textbf{Data (sequences of random variables and observed values):} For all $i\in \mathbb{N}$, denote the random variable for the $i$-th data sample as $Z_i:=(X_i, Y_i)\in \mathcal{X}\times\mathcal{Y}=\mathcal{Z}$, and denote observed value of that random variable in lower case as $z_i:=(x_i, y_i)\in \mathcal{X}\times\mathcal{Y}=\mathcal{Z}$. Similarly, denote a sequence of random variables as $Z_{1:m}:=(Z_1, ..., Z_m)$ and a sequence of observed data values as $z_{1:m}:=(z_1, ..., z_m)$ for any $m \in \mathbb{N}$. To concisely denote indices, let $[m]:=\{1, ..., m\}$.
    \item \textbf{Hyperparameter space:} Let $\Lambda \subseteq \mathbb{R}$ be the space of hyperparameters. So, $\lambda\in \Lambda$ denotes a particular hyperparameter in this space.
    \item \textbf{Loss functions:} For some real-valued function $g :\mathcal{Z}\times \Lambda \rightarrow \mathbb{R}$, and for all $\lambda\in\Lambda$ and all $i\in \mathbb{N}$, let $L_i(\lambda):= g(Z_i,\lambda)$ denote the loss function determined by a random variable $Z_i$ and parameter $\lambda\in \Lambda$, and let $\ell_i(\lambda):= g(z_i,\lambda)$ denote the loss function determined by the observed value $z_i$ and parameter $\lambda\in \Lambda$. To emphasize when no particular $\lambda$ has yet been given as an argument, we will use $L_i(\cdot):=g(Z_i, \cdot)$ to denote the loss function with randomness over $Z_i$ not yet determined, and we will use $\ell_i(\cdot):= g(z_i, \cdot)$ to denote the loss function for observed value $z_i$. 
    \item \textbf{Bags:} For any sequence $v_{1:m}:=(v_1, ..., v_m)$, let $\lbag v_{1:m} \rbag$ denote the \textit{bag} (or multiset) containing the values of that sequence $v_{1:m}$ (but importantly, $\lbag v_{1:m} \rbag$ does \textit{not} contain the information about the order of $v_{1:m}$). For example, $\lbag (1, 7, 42, 7) \rbag$ conveys that there is one 1, two 7s, and one 42 in the bag, but it does \textit{not} convey the order that these appeared in the sequence $(1, 7, 42, 7)$. In particular, with this notation, for the random variables $Z_{1:m}$, let $\lbag Z_{1:m} \rbag$ denote the bag containing those random variables; for observed values $z_{1:m}$, let $\lbag z_{1:m} \rbag$ denote the bag containing those observed values; and thus $\lbag Z_{1:m} \rbag = \lbag z_{1:m} \rbag$ denotes the event of observing $\lbag z_{1:m} \rbag$, in other words, that each $Z_i$ takes a different value in $\lbag z_{1:m} \rbag$, but it is not yet known \textit{which} value (for all $i\in [m]$).
    \item \textbf{Distributions:} Let $F_{Z_{1:m}}:=F_{Z_1, ..., Z_m}$ denote the joint distribution function for $Z_1, ..., Z_m$, and let $F_{\lbag Z_{1:m}\rbag}:=F_{\lbag Z_1, ..., Z_m\rbag}$ denote the (pushforward) distribution induced by the mapping $(z_1, ..., z_m)\rightarrow \lbag z_{1:m} \rbag$. 
\end{itemize}

% \begin{align*}
%     \hat{\lambda}_+ = \hat{\lambda}_+(\lbag L_{1:n}, B\rbag, \alpha) := \inf \Big\{\lambda_0 \in \Lambda \quad : \quad \forall \qquad \lambda \geq \lambda_0, \qquad \sum_{i=1}^{n}\frac{L_i(\lambda)}{n+1} + \frac{B}{n+1} \leq \alpha \Big\}.
% \end{align*}


% \begin{theorem}
%     Assume that $L_i(\lambda)$ is $K$-Lipschitz continuous in $\lambda$. Define  $\lambda_{\max}:=\sup\Lambda\in \Lambda$ and assume
%     \begin{align*}
%         L_i(\lambda_{\max})\leq \alpha, \qquad \sup_{\lambda}L_i(\lambda)\leq B < \infty \quad \text{almost surely}.
%     \end{align*}
%     Then, if the Lipschitz constant $K$ is known and $\hat\alpha$ exists, 
%     \begin{align*}
%         \mathbb{E}[L_{n+1}\big(\hat{\lambda}_{+}(\hat\alpha)\big)] \leq \alpha,
%     \end{align*}
%     and if $K$ is not known but $\hat\lambda_+$ is $\epsilon$-replace-one stable, then
%     \begin{align*}
%         \mathbb{E}[L_{n+1}(\hat{\lambda}_{+})] \leq \alpha + K\epsilon.
%     \end{align*}
    
    
% \end{theorem}

\begin{theorem}[Restatement of Theorem \ref{thm:gcrc_nonmonotonic}]
    Assume exchangeable $L_i(\lambda)$. Define  $\lambda_{\max}:=\sup\Lambda\in \Lambda$ and assume 
    \begin{align*}
        L_i(\lambda_{\max})\leq \alpha, \qquad \sup_{\lambda}L_i(\lambda)\leq B < \infty \quad \text{almost surely}.
    \end{align*}
    If the $L_i(\lambda)$ are $K$-Lipschitz in $\lambda$ and $\hat\lambda_+$ is $\epsilon$-replace-one stable, then
    % Then, if the Lipschitz constant $K$ is known and $\hat\alpha$ exists, 
    % \begin{align*}
    %     \mathbb{E}\big[L_{n+1}\big(\hat{\lambda}_{+}(\hat\alpha)\big)\big] \leq \alpha,
    % \end{align*}
    % and if $K$ is not known but $\hat\lambda_+$ is $\epsilon$-replace-one stable, then
    \begin{align*}
        \mathbb{E}\big[L_{n+1}(\hat\lambda_+)\big] \leq \alpha + K\epsilon.
    \end{align*}
\end{theorem}

\begin{proof} \qquad

\textbf{Step 0: Law of total expectation over draw of bag.}

Expanding $\mathbb{E}[L_{n+1}(\hat{\lambda}_{+})]$ via the law of total expectation over the event $\lbag Z_{1:n+1} \rbag = \lbag z_{1:n+1} \rbag$, we have
\begin{align}\label{eq:gcrc_exchangeable_step_0}
    \mathbb{E}\big[L_{n+1}(\hat{\lambda}_{+})\big] = \mathbb{E}_{F_{\lbag Z_{1:n+1}\rbag}}\Big[\mathbb{E}\Big[L_{n+1}(\hat{\lambda}_{+}) \mid \lbag Z_{1:n+1} \rbag = \lbag z_{1:n+1} \rbag\Big]\Big].
\end{align}
We can thus see that, to prove the desired result, it is sufficient to upper bound $\mathbb{E}\big[L_{n+1}(\hat{\lambda}_{+}) \mid \lbag Z_{1:n+1} \rbag = \lbag z_{1:n+1} \rbag\big]$ almost surely given any draw of the bag, $\lbag Z_{1:n+1}\rbag\sim F_{\lbag Z_{1:n+1}\rbag}$. 

\vspace{1cm}

\textbf{Step 1: Deterministic inequalities on bag-conditional risks.} 


Define the following family of thresholds:
\begin{align*}
    \hat{\lambda}_+(\lbag L_{1:m} \rbag, \gamma) = \inf \Big\{\lambda_0 \in \Lambda \quad : \quad \forall \qquad \text{fixed } \lambda \geq \lambda_0, \qquad \sum_{i=1}^{m}\frac{l_{i}(\lambda)}{m} \leq \gamma \Big\}.
\end{align*}
When the set is empty, define $\hat{\lambda}_+(\lbag L_{1:m} \rbag, \gamma)=\lambda_{\max}$.

By the right-continuity of $L_i(\cdot)$ (for all $i\in \mathbb{N}$), the sum $\sum_{i=1}^m\frac{L_m(\lambda)}{m}$ is right-continuous for any $Z_{1:m}$, which gives us the following deterministic inequality for any $\gamma$:
\begin{align}
    \label{eq:gcrc_exchangeable_step_1}
    \sum_{i=1}^m\frac{L_i(\lambda)}{m} \leq \gamma \qquad \forall \qquad \text{fixed } \lambda \geq \hat{\lambda}_{+}(\lbag L_{1:m}\rbag; \gamma).
\end{align}


\vspace{1cm}

\textbf{Step 2: Showing an oracle procedure deterministically controls the bag-conditional risk.}

Using the family of thresholds we defined above and applying it to the bag of loss functions for both the calibration set and the test point, define the oracle threshold as
\begin{align}\label{eq:oracle_lambda_def}
    \lambda_+^* := \hat{\lambda}_+(\lbag L_{1:n+1} \rbag, \alpha).
\end{align}

Now, consider the quantity $\mathbb{E}\big[L_{n+1}(\lambda_{+}^*) \mid \lbag Z_{1:n+1} \rbag = \lbag z_{1:n+1} \rbag\big]$, which we can view as the bag-conditional oracle risk. That is, we can interpret this quantity as the expected value that the random variable $L_{n+1}(\lambda_{+}^*)$ will take on, given the bag of data $\lbag z_{1:n+1} \rbag$. Note that conditioning on the bag of data, $\lbag z_{1:n+1} \rbag$, implies the event of observing the bag of loss functions, $\lbag l_{1:n+1} \rbag$,\footnote{That is, implies the event $\lbag L_{1}(\cdot), ..., L_{n+1}(\cdot)\rbag = \lbag l_1(\cdot), ..., l_{n+1}(\cdot) \rbag$, because the loss functions are defined as $l_i(\cdot):=g(z_i, \cdot)$ for all $i\in \mathbb{N}$.} which thereby determines the value of $\lambda_{+}^*$, since $\lambda_{+}^*$ is a symmetric function of $\lbag z_{1:n+1}\rbag$. 

This fact, that $\lambda_{+}^*$ is a symmetric function of (and thus determined by) the bag of data, implies that $L_{n+1}(\lambda_{+}^*)$ will take a value in $\lbag l_1(\lambda_{+}^*), ..., l_{n+1}(\lambda_{+}^*) \rbag$ and the specific value $L_{n+1}(\lambda_{+}^*)$ will take on is determined by the draw of $Z_{n+1}$ from $\lbag z_{1:n+1} \rbag$; that is, given $\lbag z_{1:n+1} \rbag$, we have the equivalency $L_{n+1}(\lambda_{+}^*) = l_i(\lambda_{+}^*)  \iff Z_{n+1} = z_i$ due to $g(\cdot,\lambda_{+}^*)$ inducing a bijection from $\lbag z_1, ..., z_{n+1} \rbag$ to the bag of loss values $\lbag l_{1}(\lambda_{+}^*), ..., l_{n+1}(\lambda_{+}^*) \rbag$.

First using this observation and the definition of conditional expectation, and secondly using the definition of conditional probability, we have
\begin{align*}
    \mathbb{E}\big[L_{n+1}(\lambda_{+}^*) \mid \lbag Z_{1:n+1} \rbag = \lbag z_{1:n+1} \rbag\big] & = \sum_{i=1}^{n+1}l_i(\lambda_{+}^*)\cdot \mathbb{P}\big(Z_{n+1} = z_{i} \mid \lbag Z_{1:n+1} \rbag = \lbag z_{1:n+1} \rbag \big) \\
    & = \sum_{i=1}^{n+1}l_i(\lambda_{+}^*)\cdot \frac{\mathbb{P}\big(Z_{n+1}=z_{i}, \ \lbag Z_{1:n+1} \rbag = \lbag z_{1:n+1} \rbag \big)}{\mathbb{P}\big(\lbag Z_{1:n+1} \rbag = \lbag z_{1:n+1} \rbag\big)}.
\end{align*}

Applying the law of total probability over permutations ($\sigma:[n+1]\rightarrow[n+1]$) to the probabilities in the numerator and denominator of the RHS, and then simplifying using exchangeability, we have
\begin{align*}
    \mathbb{E}\big[L_{n+1}(\lambda_{+}^*) \mid \lbag Z_{1:n+1} \rbag =  \lbag z_{1:n+1} \rbag\big]
    & = \sum_{i=1}^{n+1}l_i(\lambda_{+}^*)\cdot \frac{\sum_{\sigma:\sigma(n+1)=i}f(z_{\sigma(1)}, ..., z_{\sigma(n+1)})}{\sum_{\sigma}f(z_{\sigma(1)}, ..., z_{\sigma(n+1)})} \\
    & = \sum_{i=1}^{n+1}l_i(\lambda_{+}^*)\cdot \frac{\sum_{\sigma:\sigma(n+1)=i}f(z_1, ..., z_{n+1})}{\sum_{\sigma}f(z_1, ..., z_{n+1})} \\
    & = \sum_{i=1}^{n+1}l_i(\lambda_{+}^*)\cdot \frac{n!}{(n+1)!} \\
    & = \sum_{i=1}^{n+1}\frac{l_i(\lambda_{+}^*)}{n+1}
\end{align*}

Now, using Eq. \eqref{eq:gcrc_exchangeable_step_1}, 
\begin{align}\label{eq:gcrc_exchangeable_step_2}
    &\sum_{i=1}^{n+1}\frac{l_{i}(\lambda)}{n+1}\leq \alpha \qquad \forall \qquad \text{fixed } \lambda \geq \lambda_{+}^*,
\end{align}
where here by ``fixed'' we mean that $\lambda$ is constant conditional on the bag of data, $\lbag Z_{1:n+1} \rbag = \lbag z_{1:n+1} \rbag$.

\vspace{1cm}

\textbf{Step 3: Relating the empirically-selected hyperparameter to the oracle parameter}


In Step 2, we have shown that an oracle procedure deterministically controls the risk, conditional on the event $\lbag Z_{1:n+1} \rbag = \lbag z_{1:n+1} \rbag$. However, this oracle procedure cannot be implemented in practice due to requiring knowledge of the test point loss, $L_{n+1}$; that is, recall we defined $\lambda_+^* := \hat{\lambda}_+(\lbag L_{1:n+1} \rbag, \alpha)$. In this last step of the proof we would now like to prove a deterministic relationship between the oracle parameter, $\lambda_+^*$, and the empirically-selected parameter, $\hat{\lambda}_+:=\hat{\lambda}_+(\lbag L_{1:n}, B \rbag, \alpha)$, to conclude valid risk control for the empirical method; that is, we would like to deterministically show $\lambda_+^* \leq \hat{\lambda}_+$.


To begin, recall that in Step 2, conditioning on the event $\lbag Z_{1:n+1} \rbag = \lbag z_{1:n+1} \rbag$ implied that $\hat{\lambda}_+(\lbag L_{1:n+1} \rbag, \alpha) = \hat{\lambda}_+(\lbag l_{1:n+1} \rbag, \alpha)$. For the empirical parameter, however, conditioning on $\lbag Z_{1:n+1} \rbag = \lbag z_{1:n+1} \rbag$ does \textit{not} imply what may seem analogous: that is, it could be that $\hat{\lambda}_+(\lbag L_{1:n}, B \rbag, \alpha) \neq \hat{\lambda}_+(\lbag l_{1:n}, B \rbag, \alpha)$. This is because the event $\lbag Z_{1:n+1} \rbag = \lbag z_{1:n+1} \rbag$ only allows us to conclude that the $n$ random variables in $\lbag L_{1:n}\rbag$ take $n$ distinct values from the bag of $n+1$ observed losses, $\lbag l_{1:n+1}\rbag$, but we do not yet know \textit{which} observed losses. That is, what we \textit{can} conclude is, given $\lbag Z_{1:n+1} \rbag = \lbag z_{1:n+1} \rbag$, that
\begin{align*}
    \lbag Z_{1:n+1} \rbag = \lbag z_{1:n+1} \rbag \implies \lbag L_{1:n} \rbag \in \Big\{\lbag l_{1:n+1\backslash i} \rbag \Big\}_{i\in [n+1]},
\end{align*}
where $\lbag l_{1:n+1\backslash i} \rbag := \lbag l_{1:n+1} \rbag \backslash \{l_i\}$. Accordingly, conditioning on the event $\lbag Z_{1:n+1} \rbag = \lbag z_{1:n+1} \rbag$ implies the following equivalence:
\begin{align}\label{eq:oracle_empirical_inequality_equivalences}
    \lambda_+^* \leq \hat\lambda_+(\lbag L_{1:n}, B \rbag, \alpha) \ \text{almost surely} \ & \iff \lambda_+^* \leq \hat\lambda_+(\lbag l_{1:n+1\backslash i}, B \rbag, \alpha) \quad \forall \quad i \in [n+1].
\end{align}
So, we will focus on proving the latter. 




Recall that we have assumed that $l$ is bounded above by $B$ (i.e., $\sup_{\lambda}l_i(\lambda)\leq B < \infty$), for any $\lambda\in \Lambda$, deterministically we have
\begin{align}\label{eq:step_3_conservative_adjustment}
    \sum_{j\in[n+1]}\frac{l_j(\lambda)}{n+1}
    \quad \leq \quad  \frac{B}{n+1} + \sum_{j\in[n+1]\backslash i}\frac{l_j(\lambda)}{n+1} \qquad \forall \qquad i \in [n+1],
\end{align}


where the right-hand side of the above inequalities can be interpreted as a conservative risk obtained by replacing $l_i(\lambda)$ with the bound $B$.


Thus, for any $\lambda$, we can use \eqref{eq:step_3_conservative_adjustment} to conclude
\begin{align}\label{eq:step_3_risk_relation}
    \frac{B}{n+1} + \sum_{j\in[n+1]\backslash i}\frac{l_j(\lambda)}{n+1} \leq \alpha \implies & \sum_{j\in[n+1]}\frac{l_j(\lambda)}{n+1} \leq \alpha \qquad \forall \qquad i \in [n+1].
\end{align}



So, by \eqref{eq:step_3_risk_relation} and the definition of $\hat{\lambda}_+(\lbag l_{1:n+1\backslash i}, B \rbag)$, it must hold that 
\begin{align}\label{eq:pf_step_3_result}
    \lambda_+^* \leq & \ \hat{\lambda}_+(\lbag l_{1:n+1\backslash i}, B \rbag) \quad \forall \quad i \in [n+1].
\end{align}

\vspace{1cm}

\textbf{Step 4: Analyzing the bag-conditional risk of the empirical procedure via Lipschitz continuity}

Now we turn to plugging the empirical algorithm's hyperparameter into the bag-conditional risk, and then analyzing this risk using Lipschitz continuity and a type of stability called ``replace-one'' stability. This analysis is necessary because $\hat\lambda_+=\widehat{\lambda}_+(\lbag Z_{1:n}, B\rbag, \alpha)$ is not constant conditional on $\lbag Z_{1:n+1} \rbag = \lbag z_{1:n+1} \rbag$, so Eq. \eqref{eq:gcrc_exchangeable_step_2} would not direclty apply. 

Begin by expanding this quantity by LOTE over the event $Z_{n+1}=z_i$:
\begin{align}\label{eq:empirical_bag_cond_LOTE}
    \mathbb{E}\Big[&L_{n+1}(\hat\lambda_{+}) \mid \lbag Z_{1:n+1} \rbag = \lbag z_{1:n+1} \rbag\Big] \nonumber \\
    & = \mathbb{E}\Big[\mathbb{E}\Big[L_{n+1}(\hat\lambda_{+}) \mid \lbag Z_{1:n+1} \rbag = \lbag z_{1:n+1} \rbag, Z_{n+1}=z_i\Big] \mid \lbag Z_{1:n+1} \rbag = \lbag z_{1:n+1} \rbag\Big] \nonumber \\
    & = \sum_{i=1}^{n+1}\mathbb{E}\Big[L_{n+1}(\hat\lambda_{+}) \mid \lbag Z_{1:n+1} \rbag = \lbag z_{1:n+1} \rbag, Z_{n+1}=z_i\Big]\cdot \mathbb{P}\big(Z_{n+1}=z_i \mid \lbag Z_{1:n+1} \rbag = \lbag z_{1:n+1} \rbag\big).
\end{align}

Now, note that we can write the inner expectation as
\begin{align*}
    \mathbb{E}\Big[L_{n+1}(\hat\lambda_{+}) \mid \lbag Z_{1:n+1} \rbag = \lbag z_{1:n+1} \rbag, Z_{n+1}=z_i\Big] & = \mathbb{E}\Big[L_{n+1}(\hat\lambda_{+}) \mid \lbag Z_{1:n} \rbag = \lbag z_{1:n+1\backslash i} \rbag, Z_{n+1}=z_i\Big] \\
    & = \mathbb{E}\Big[L_{n+1}(\hat\lambda_{+}) \mid \lbag L_{1:n} \rbag = \lbag l_{1:n+1\backslash i} \rbag, L_{n+1}=l_i\Big].
\end{align*}

And recalling that we defined $\hat\lambda_+:= \widehat{\lambda}_+\big(\lbag L_{1:n}, B \rbag, \ \alpha\big)$, we have
\begin{align*}
    \mathbb{E}\Big[&L_{n+1}(\hat\lambda_{+}) \mid \lbag Z_{1:n+1} \rbag = \lbag z_{1:n+1} \rbag, Z_{n+1}=z_i\Big]  \\
    & = \mathbb{E}\Big[L_{n+1}\Big(\widehat{\lambda}_+\big(\lbag L_{1:n}, B  \rbag, \ \alpha\big)\Big) \mid \lbag L_{1:n} \rbag = \lbag l_{1:n+1\backslash i} \rbag, L_{n+1}=l_i\Big] \\
    & = \mathbb{E}\Big[l_{i}\Big(\widehat{\lambda}_+\big(\lbag l_{1:n+1\backslash i}, B  \rbag, \ \alpha\big)\Big) \mid \lbag L_{1:n} \rbag = \lbag l_{1:n+1\backslash i} \rbag, L_{n+1}=l_i\Big] \\
    & = l_{i}\Big(\widehat{\lambda}_+\big(\lbag l_{1:n+1\backslash i}, B  \rbag, \ \alpha\big)\Big).
\end{align*}


So, plugging this into Eq. \eqref{eq:empirical_bag_cond_LOTE}, we have
\begin{align*}
    \mathbb{E}\Big[& L_{n+1}(\hat\lambda_{+}) \mid \lbag Z_{1:n+1} \rbag = \lbag z_{1:n+1} \rbag\Big] \\
    & = \sum_{i=1}^{n+1}l_{i}\Big(\widehat{\lambda}_+\big(\lbag l_{1:n+1\backslash i}, B  \rbag, \ \alpha\big)\Big) \cdot \mathbb{P}\big(Z_{n+1}=z_i \mid \lbag Z_{1:n+1} \rbag = \lbag z_{1:n+1} \rbag\big) \\
    & = \sum_{i=1}^{n+1}l_{i}\Big(\widehat{\lambda}_+\big(\lbag l_{1:n+1\backslash i}, B  \rbag, \ \alpha\big)\Big) \cdot\frac{\mathbb{P}\big(Z_{n+1}=z_{i}, \ \lbag Z_{1:n+1} \rbag = \lbag z_{1:n+1} \rbag \big)}{\mathbb{P}\big(\lbag Z_{1:n+1} \rbag = \lbag z_{1:n+1} \rbag\big)}
\end{align*}


Applying the law of total probability over permutations ($\sigma:[n+1]\rightarrow[n+1]$) to the probabilities in the numerator and denominator of the RHS, and then simplifying using exchangeability, we have
\begin{align}\label{eq:gcrc_empirical_risk_bag_conditional}
    \mathbb{E}\big[L_{n+1}(\hat\lambda_{+}) \mid \lbag Z_{1:n+1} \rbag =  \lbag z_{1:n+1} \rbag\big] \nonumber
    & = \sum_{i=1}^{n+1}l_{i}\Big(\widehat{\lambda}_+\big(\lbag l_{1:n+1\backslash i}, B  \rbag, \ \alpha\big)\Big)\cdot \frac{\sum_{\sigma:\sigma(n+1)=i}f(z_{\sigma(1)}, ..., z_{\sigma(n+1)})}{\sum_{\sigma}f(z_{\sigma(1)}, ..., z_{\sigma(n+1)})} \nonumber \\
    & = \sum_{i=1}^{n+1}l_{i}\Big(\widehat{\lambda}_+\big(\lbag l_{1:n+1\backslash i}, B  \rbag, \ \alpha\big)\Big)\cdot \frac{\sum_{\sigma:\sigma(n+1)=i}f(z_1, ..., z_{n+1})}{\sum_{\sigma}f(z_1, ..., z_{n+1})} \nonumber \\
    & = \sum_{i=1}^{n+1}l_{i}\Big(\widehat{\lambda}_+\big(\lbag l_{1:n+1\backslash i}, B  \rbag, \ \alpha\big)\Big)\cdot \frac{n!}{(n+1)!} \nonumber \\
    & = \frac{1}{n+1}\sum_{i=1}^{n+1}l_{i}\Big(\widehat{\lambda}_+\big(\lbag l_{1:n+1\backslash i}, B  \rbag, \ \alpha\big)\Big).
\end{align}

Importantly, note that in Eq. 
\eqref{eq:gcrc_empirical_risk_bag_conditional}, which is the true bag-conditional risk of the empirical algorithm, that the value of $\widehat{\lambda}_+\big(\lbag l_{1:n+1\backslash i}, B  \rbag, \ \alpha\big)$ differs for each loss function, $l_i$, in the summation, which is why we cannot apply Eq. \eqref{eq:gcrc_exchangeable_step_2} to bound it directly. Instead, we relate Eq. \eqref{eq:gcrc_empirical_risk_bag_conditional} and Eq. \eqref{eq:gcrc_exchangeable_step_2} using Lipschitz continuity and the replace-one (in)stability of the algorithm. 



For all $i\in[n]$, denote the magnitude of perturbation to the hyperparameter $\hat\lambda_+(\lbag l_{1:n}, B\rbag, \alpha)$ caused by replacing the loss function $l_i$ with $l_{n+1}$ as 
\begin{align*}
    \epsilon_i = \Big|\hat\lambda_+(\lbag l_{1:n}, B\rbag, \alpha)-\hat\lambda_+(\lbag l_{1:n+1\backslash i}, B\rbag, \alpha)\Big|.
\end{align*}
Then, by $K$-Lipschitz continuity, we know that for all $i\in [n]$, that
\begin{align*}
    \Big|l_i\big(\hat\lambda_+(\lbag l_{1:n}, B\rbag, \alpha)\big) - l_i\big(\hat\lambda_+(\lbag l_{1:n+1\backslash i}, B\rbag, \alpha)\big)\Big| \leq K \epsilon_i.
\end{align*}

Then, note that 
\begin{align*}
    \sum_{i=1}^{n+1}\frac{l_i\big(\hat\lambda_+(\lbag l_{1:n+1\backslash i}, B\rbag, \alpha)\big)}{n+1} & \leq \sum_{i=1}^{n+1}\frac{l_i\big(\hat\lambda_+(\lbag l_{1:n}, B\rbag, \alpha)\big) + K\epsilon_i}{n+1} = \sum_{i=1}^{n+1}\frac{l_i\big(\hat\lambda_+(\lbag l_{1:n}, B\rbag, \alpha)\big)}{n+1} + \frac{K}{n+1}\sum_{i=1}^{n+1}\epsilon_i 
\end{align*}

By Eq. \eqref{eq:pf_step_3_result}, we know that $\hat\lambda_+(\lbag l_{1:n}, B\rbag, \alpha) \geq \lambda_+^*$, and because this $\hat\lambda_+(\lbag l_{1:n}, B\rbag, \alpha)$ does not depend on $i$, we can apply Eq. \eqref{eq:gcrc_exchangeable_step_2} to bound it above by $\alpha$, which gives 
\begin{align*}
     \sum_{i=1}^{n+1}\frac{l_i\big(\hat\lambda_+(\lbag l_{1:n+1\backslash i}, B\rbag, \alpha)\big)}{n+1} & \leq \alpha + \frac{K}{n+1}\sum_{i=1}^{n+1}\epsilon_i \\
     \iff \mathbb{E}\big[L_{n+1}(\hat\lambda_{+}) \mid \lbag Z_{1:n+1} \rbag =  \lbag z_{1:n+1} \rbag\big] & \leq \alpha + \frac{K}{n+1}\sum_{i=1}^{n+1}\epsilon_i.
\end{align*}
Then, marginalizing over the event and applying $\epsilon$-replace-one stability, we have
\begin{align*}
    \mathbb{E}\big[L_{n+1}(\hat\lambda_{+}) \big] & \leq \alpha + K \epsilon.
\end{align*}

% Moreover, if the Lipschitz constant $K$ is known, then we can find a conservative empirical significance level, $\hat\alpha$, that can be implemented to control the risk at the target level $\alpha$. To do so, define
% \begin{align*}
%     \hat\epsilon_i = \max_{b\in \{0, B\}}\Big|\hat\lambda_+(\lbag l_{1:n}, B\rbag, \alpha)-\hat\lambda_+(\lbag l_{1:n\backslash i}, b, B\rbag, \alpha)\Big|,
% \end{align*}
% which is a conservative estimate of $\epsilon_i$ using only the observations $l_{1:n}$, as note that by construction, $\hat\epsilon_i\geq \epsilon_i$. Then, fit the conservative significance level as 
% \begin{align*}
%     \hat\alpha := \sup\Big(\Big\{\alpha' \leq \alpha \quad : \quad \alpha' + \frac{K}{n+1}\sum_{i=1}^{n+1}\hat\epsilon_i(\alpha')\leq \alpha\Big\}  \Big), 
% \end{align*}
% and by a similar argument as above, if $\hat\alpha$ exists then 
% \begin{align*}
%     \mathbb{E}\big[L_{n+1}\big(\hat\lambda_{+}(\hat\alpha)\big) \big] & \leq \alpha.
% \end{align*}

Note that the above analysis also applies for any alternative $\hat\lambda_+'$ that is a deterministic function of $\lbag l_{1:n}\rbag$ and where almost surely $\hat\lambda_+'\geq \lambda_+^*$, and assuming, $\hat\lambda_+'$ is $\epsilon$-replace-one stable. 

    
\end{proof}



\subsection{Proof for Conformal Policy Control Guarantee}
\label{app:cpc_proof}




\textbf{Notation:}
\begin{itemize}
    \item \textbf{Data (sequences of random variables and observed values):} For all $i\in \mathbb{N}$, denote the random variable for the $i$-th data sample as $Z_i\in \mathcal{Z}$, and denote observed value of that random variable in lower case as $z_i\in\mathcal{Z}$. For example, we might have $\mathcal{Z}=\mathcal{X}\times \mathcal{A}\times \mathbb{R}\times \mathbb{R}$ if $Z=(X, A, R, L)$ is a vector of context, action, reward, and loss. Similarly, denote a sequence of random variables as $Z_{1:m}:=(Z_1, ..., Z_m)$ and a sequence of observed data values as $z_{1:m}:=(z_1, ..., z_m)$ for any $m \in \mathbb{N}$. To concisely denote indices, let $[m]:=\{1, ..., m\}$.
    \item \textbf{Hyperparameters:} Let $\mathbb{R}$ be the space of hyperparameters, and let $\beta\in \mathbb{R}$ denote a particular hyperparameter. 
    \item \textbf{Loss functions:} For some real-valued function $g :\mathcal{Z}\times \beta \rightarrow \mathbb{R}$, and for all $\beta\in\mathbb{R}$ and all $i\in \mathbb{N}$, let $L_i(\beta):= g(Z_i,\beta)$ denote the loss function determined by a random variable $Z_i$ and parameter $\beta\in \mathbb{R}$, and let $\ell_i:= g(z_i,\beta)$ denote the loss function determined by the observed value $z_i$ and parameter $\beta\in \mathbb{R}$. Note that allowing the loss functions to be a function of $\beta$ does not preclude the setting where the losses are independent of $\beta$, that is where $L_i(\beta)=L_i$ for all $\beta\in\mathbb{R}$ (i.e., where changing $\beta$ does not change the loss). Such a setting motivates allowing $\beta$ to parameterize the policies.
    \item \textbf{Bags:} For any sequence $v_{1:m}:=(v_1, ..., v_m)$, let $\lbag v_{1:m} \rbag$ denote the \textit{bag} (or multiset) containing the values of that sequence $v_{1:m}$ (but importantly, $\lbag v_{1:m} \rbag$ does \textit{not} contain the information about the order of $v_{1:m}$). For example, $\lbag (1, 7, 42, 7) \rbag$ conveys that there is one 1, two 7s, and one 42 in the bag, but it does \textit{not} convey the order that these appeared in the sequence $(1, 7, 42, 7)$. In particular, with this notation, for the random variables $Z_{1:m}$, let $\lbag Z_{1:m} \rbag$ denote the bag containing those random variables; for observed values $z_{1:m}$, let $\lbag z_{1:m} \rbag$ denote the bag containing those observed values; and thus $\lbag Z_{1:m} \rbag = \lbag z_{1:m} \rbag$ denotes the event of observing $\lbag z_{1:m} \rbag$, in other words, that each $Z_i$ takes a different value in $\lbag z_{1:m} \rbag$, but it is not yet known \textit{which} value (for all $i\in [m]$).
    \item \textbf{Policies (Distributions):} Let $\pi_{1:m}:=\pi_{Z_1, ..., Z_m}$ denote the policy or joint density function\footnote{In a slight abuse of notation for the policy control setting, we do not distinguish between CDFs and density functions here.} for $Z_1, ..., Z_m$, and let $\pi_{\lbag Z_{1:m}\rbag}:=\pi_{\lbag Z_1, ..., Z_m\rbag}$ denote the (pushforward) distribution induced by the mapping $(z_0, ..., z_m)\rightarrow \lbag z_{1:m} \rbag$. 
\end{itemize}




% \begin{theorem}
% \label{appthm:cpc}
%     Assume that $L_i\cdot w_i^{(\beta)}$ is $K$-Lipschitz in $\beta$, that $\mathbb{E}_{\pi_0}[L_t]\leq \alpha$, and  $L_i\leq B < \infty$ for all $i$.
%     Then, if the Lipschitz constant $K$ is known and $\hat\alpha$ exists, 
%     \begin{align*}
%         \mathbb{E}_{\pi_{0:t}^{(\hat\beta(\hat\alpha))}}\big[L_{t} \big] & \leq \alpha.
% \end{align*}
%     and if $K$ is not known but $\hat\beta$ is $\epsilon$-replace-one stable, then
%     \begin{align*}
%         \mathbb{E}_{\pi_{0:t}^{(\hat\beta)}}\big[L_{t} \big] & \leq \alpha + K\cdot \epsilon .
% \end{align*}
% \end{theorem}

\begin{theorem}[Restatement of Theorem \ref{thm:cpc}]
% \label{appthm:cpc}
    Assume that $\mathbb{E}_{\pi_0}[L_t]\leq \alpha$,  $L_i\in[0,B], B<\infty$ for all $i$, and $\hat\beta$ is $\epsilon$-replace-one stable. If $w_i^{(\beta)}$ is $K$-Lipschitz in $\beta$ for all $i$, then
    \begin{align*}
        \mathbb{E}_{\pi_{0:t}^{(\hat\beta)}}\big[L_{t} \big] \leq \alpha + BK\epsilon.
    \end{align*}
    If $(l_i - \alpha) \cdot w_i^{(\beta)}$ is $K'$-Lipschitz in $\beta$, then
    \begin{align*}
        \mathbb{E}_{\pi_{0:t}^{(\hat\beta)}}\big[L_{t} \big] \leq \alpha + K'\epsilon.
    \end{align*}
\end{theorem}

\begin{proof} \qquad


For ease of notation, we conduct the proof in an online setting where only one action is taken (or sample is generated) at each policy improvement step. This will allow the indices $t=0, 1, ...$ to denote both the timestep of policy improvement and the timestep of the action/sample, but the result will hold generally for a batch setting too (with the guarantee still holding marginally over the whole procedure). 

We proceed by induction. For $t=0$, we assumed access to a safe initial policy, $\pi_0$, satisfying our risk control criterion. That is, we assume
\begin{align*}
    \mathbb{E}_{Z_0\sim \pi_0}[L_0] \leq \alpha.
\end{align*}
Now, for the induction hypothesis assume that risk control holds at time $t-1$, that is, assume
\begin{align*}
    \mathbb{E}_{Z_{0:t-1}\sim \pi_{0:t-1}^{(\hat\beta_1:\hat\beta_{t-1})}}[L_{t-1}] \leq \alpha.
\end{align*}

We then want to show that risk control also holds at time $t$:
\begin{align*}
    \mathbb{E}_{Z_{0:t}\sim \pi_{0:t}^{(\hat\beta_1:\hat\beta_t)}}[L_{t}] \leq \alpha,
\end{align*}
where note that $\pi_{0:t}^{(\hat\beta_t)}(z_{0:t}) = \pi_{t}^{(\hat\beta_t)}(z_t \mid z_{0:t-1})\pi_{t-1}^{(\hat\beta_1:\hat\beta_{t-1})}(z_{0:t-1})$. Note that at time $t$, we will only consider tuning $\hat\beta_t$, and we treat $\hat\beta_1, ..., \hat\beta_{t-1}$ as previously tuned. Accordingly, to lighten notation, we hereon denote the expectation we wish to bound as $\mathbb{E}_{\pi_{0:t}^{(\hat\beta)}}[L_{t}]:=\mathbb{E}_{Z_{0:t}\sim \pi_{0:t}^{(\hat\beta_1:\hat\beta_t)}}[L_{t}]$, where note that we also let $\hat\beta:=\hat\beta_t$.

For any control parameter $\beta$, any policies $\pi_{0:t}^{(\beta)}$, and any bag of data $\lbag z_{0:t}\rbag$, following \citet{tibshirani2019conformal, prinster2024conformal} define the oracle normalized conformal weights as
\begin{align}\label{eq:cp_weights_def_cpc}
    \tilde{w}_i^{(\beta)} := \tilde{w}_i\big(\pi_{0:t}^{(\beta)}, \lbag z_{0:t}\rbag\big) = \mathbb{P}_{\pi_{0:t}^{(\beta)}}(Z_t=z_i \mid \lbag z_{0:t}\rbag) =\frac{\sum_{\sigma:\sigma(t)=i}\pi_{0:t}^{(\beta)}(z_{\sigma(0)}, ..., z_{\sigma(t)})}{\sum_{\sigma}\pi_{0:t}^{(\beta^*)}(z_{\sigma(0)}, ..., z_{\sigma(t)})}
\end{align}
for all $i\in \{0, ..., t\}$, as well as the unnormalized weights,
\begin{align}\label{eq:cp_weights_def_cpc_unnormed}
    w_i^{(\beta)} := w_i\big(\pi_{0:t}^{(\beta)}, \lbag z_{0:t}\rbag\big) = \sum_{\sigma:\sigma(t)=i}\pi_{0:t}^{(\beta)}(z_{\sigma(0)}, ..., z_{\sigma(t)}),
\end{align}
so note that $\tilde{w}_i^{(\beta)} = \frac{w_i^{(\beta)}}{\sum_{i=0}^tw_i^{(\beta)}}$,
and define the vector $\tilde{w}_{0:t}^{(\beta)}:=(\tilde{w}_0^{(\beta)}, ..., \tilde{w}_t^{(\beta)})$.

Also define the conservative (normalized) weight for an unknown test point as
\begin{align*}
    \tilde{w}_{\max}^{(\beta)}:=\sup_{z_t\in \mathcal{Z}}(\tilde{w}_t^{(\beta)}).
\end{align*}

Assume there is some small, positive $\beta_{\min}\ll1$ such that $\beta_{\min}\leq \pi_t(x)/\pi_0(x)$ for all $x\in \mathcal{X}$. This is the safe hyperparameter value, for which $\pi_t^{(\beta_{\min})}=\pi_0$.

\vspace{1cm}


\textbf{Step 0: Law of total expectation over draw of bag.}

Expanding $\mathbb{E}_{\pi_{0:t}^{(\hat\beta)}}[L_{t}]$ via the law of total expectation over the event $\lbag Z_{0:t} \rbag = \lbag z_{0:t} \rbag$, we have
\begin{align}\label{eq:cpc_step_0}
    \mathbb{E}_{\pi_{0:t}^{(\hat\beta)}}\big[L_{t}\big] = \mathbb{E}\Big[\mathbb{E}_{\pi_{0:t}^{(\hat\beta)}}\Big[L_{t} \mid \lbag Z_{0:t} \rbag = \lbag z_{0:t} \rbag\Big]\Big].
\end{align}
We can thus see that, to prove the desired result, it is sufficient to upper bound $\mathbb{E}_{\pi_{0:t}^{(\hat\beta)}}\big[L_{t} \mid \lbag Z_{0:t} \rbag = \lbag z_{0:t} \rbag\big]$ almost surely given any draw of the bag, $\lbag Z_{0:t}\rbag$. 


\vspace{1cm}

\textbf{Step 1: Deterministic inequalities on bag-conditional risks.} 


Define the following family of hyperparameters:
\begin{align*}
    \hat{\beta}(\lbag l_{1:m} \rbag, \tilde{w}_{1:m}^{(\cdot)},  \gamma) = \sup \Big\{\beta' \in \beta \quad : \quad \forall \qquad \beta \leq \beta', \qquad \sum_{i=1}^ml_i\cdot \tilde{w}_i^{(\beta)} \leq \gamma \Big\}.
\end{align*}
When the set is empty, define $\hat{\beta}(\lbag l_{1:m} \rbag, \tilde{w}_{1:m}^{(\cdot)},  \gamma)=\beta_{\min}$.

By the continuity of $\tilde{w}_{i}^{(\cdot)}$ (and $L_i(\cdot)$) for all $i\in \mathbb{N}$, the sum $\sum_{i=1}^ml_i\cdot \tilde{w}_i^{(\beta)}$ is right-continuous for any $Z_{1:m}$, which gives us the following deterministic inequality for any $\gamma$:
\begin{align}
    \label{eq:cpc_step_1}
    \sum_{i=1}^ml_i\cdot \tilde{w}_i^{(\beta)} \leq \gamma \qquad \forall \qquad \beta \leq \hat{\beta}(\lbag L_{1:m} \rbag, \tilde{w}_{1:m}^{(\cdot)},  \gamma).
\end{align}


\vspace{1cm}

\textbf{Step 2: Showing an oracle procedure deterministically controls the bag-conditional risk.}

Using the family of thresholds we defined above and applying it to the bag of loss functions for both the calibration set and the test point, define the oracle threshold as
\begin{align}\label{eq:cpc_oracle_lambda_def}
    \beta^* := \hat{\beta}(\lbag L_{0:t} \rbag, \tilde{w}_{0:t}^{(\cdot)},  \alpha).
\end{align}

Now, consider the quantity $\mathbb{E}_{\pi_{0:t}^{(\beta^*)}}\big[L_{t} \mid \lbag Z_{0:t} \rbag = \lbag z_{0:t} \rbag\big]$, which we can view as the bag-conditional oracle risk. That is, we can interpret this quantity as the expected value, with respect to the policy $\pi_{0:t}^{(\beta^*)}$, that the random variable $L_{t}$ will take on, given the bag of data $\lbag z_{0:t} \rbag$. Note that conditioning on the bag of data, $\lbag z_{0:t} \rbag$, implies the event of observing the bag of loss functions, $\lbag l_{0:t} \rbag$,\footnote{That is, implies the event $\lbag L_{1}, ..., L_{t}\rbag = \lbag l_1, ..., l_{t} \rbag$, because the loss functions are defined as $l_i:=g(z_i)$ for all $i\in \mathbb{N}$.} which thereby determines the value of $\beta^*$, since $\beta^*$ is a function of $\lbag z_{0:t}\rbag$ (as well as the policy, $\pi_{0:t}^{(\cdot)}$, and the constraint function $g$). 

Using the definition of conditional expectation, and secondly using the definition of conditional probability, we have
\begin{align*}
    \mathbb{E}_{\pi_{0:t}^{(\beta^*)}}\big[L_{t} \mid \lbag Z_{0:t} \rbag = \lbag z_{0:t} \rbag\big] & = \sum_{i=0}^{t}l_i\cdot \mathbb{P}\big(Z_{t} = z_{i} \mid \lbag Z_{0:t} \rbag = \lbag z_{0:t} \rbag \big) \\
    & = \sum_{i=0}^{t}l_i\cdot \frac{\mathbb{P}\big(Z_{t}=z_{i}, \ \lbag Z_{0:t} \rbag = \lbag z_{0:t} \rbag \big)}{\mathbb{P}\big(\lbag Z_{0:t} \rbag = \lbag z_{0:t} \rbag\big)}.
\end{align*}

Applying the law of total probability over permutations ($\sigma:\{0, ..., t\}\rightarrow\{0, ..., t\}$) to the probabilities in the numerator and denominator of the RHS, we have
\begin{align*}
    \mathbb{E}_{\pi_{0:t}^{(\beta^*)}}\big[L_{t} \mid \lbag Z_{0:t} \rbag =  \lbag z_{0:t} \rbag\big]
    & = \sum_{i=0}^{t}l_i\cdot \frac{\sum_{\sigma:\sigma(t)=i}\pi_{0:t}^{(\beta^*)}(z_{\sigma(0)}, ..., z_{\sigma(t)})}{\sum_{\sigma}\pi_{0:t}^{(\beta^*)}(z_{\sigma(0)}, ..., z_{\sigma(t)})} \\
    & = \sum_{i=0}^{t}l_i\cdot \tilde{w}_i^{(\beta^*)}
\end{align*}

Now, using Eq. \eqref{eq:cpc_step_1}, 
\begin{align}\label{eq:cpc_step_2}
    &\sum_{i=0}^{t}l_i\cdot \tilde{w}_i^{(\beta)}\leq \alpha \qquad \forall \qquad \text{fixed } \beta \leq \beta^*,
\end{align}
where here by ``fixed'' we mean that $\beta$ is constant conditional on the bag of data, $\lbag Z_{0:t} \rbag = \lbag z_{0:t} \rbag$.

\vspace{1cm}


\textbf{Step 3: Relating the empirically-selected hyperparameter to the oracle parameter}


In Step 2, we have shown that an oracle procedure deterministically controls the risk, conditional on the event $\lbag Z_{0:t} \rbag = \lbag z_{0:t} \rbag$. However, this oracle procedure cannot be implemented in practice due to requiring knowledge of the test point loss, $L_{t}$, as well as the weight assigned to the test point (as a function of $\beta$), $\tilde{w}_{t}^{\cdot}$; that is, recall we defined $\beta^* := \hat{\beta}(\lbag L_{0:t} \rbag, \tilde{w}_{0:t}^{(\cdot)},  \alpha)$. In this next step of the proof we would now like to prove a deterministic relationship between the oracle parameter, $\beta^*$, and the empirically-selected parameter, $\hat{\beta}:=\hat{\beta}(\lbag L_{1:t-1}, B \rbag, \tilde{w}_{0:t-1}^{(\cdot)}, \tilde{w}_{\max}^{(\cdot)},  \alpha)$, to conclude valid risk control for the empirical method; that is, we would like to deterministically show $\beta^* \geq \hat{\beta}$.


To begin, recall that in Step 2, conditioning on the event $\lbag Z_{0:t} \rbag = \lbag z_{0:t} \rbag$ implied that $\hat{\beta}(\lbag L_{0:t} \rbag, \tilde{w}_{0:t}^{(\cdot)},  \alpha) = \hat{\beta}(\lbag l_{0:t} \rbag, \tilde{w}_{0:t}^{(\cdot)},  \alpha)$. For the empirical parameter, however, conditioning on $\lbag Z_{0:t} \rbag = \lbag z_{0:t} \rbag$ does \textit{not} imply what may seem analogous: that is, it could be that $\hat{\beta}(\lbag L_{0:t-1}, B \rbag, \tilde{w}_{0:t-1}^{(\cdot)}, \tilde{w}_{\max}^{(\cdot)}, \alpha) \neq \hat{\beta}(\lbag l_{0:t-1}, B \rbag, \tilde{w}_{0:t-1}^{(\cdot)}, \tilde{w}_{\max}^{(\cdot)}, \alpha)$. This is because the event $\lbag Z_{0:t} \rbag = \lbag z_{0:t} \rbag$ only allows us to conclude that the $t-1$ random variables in $\lbag L_{1:t-1}\rbag$ take $t-1$ distinct values from the bag of $t$ observed losses, $\lbag l_{0:t}\rbag$, but we do not yet know \textit{which} observed losses. That is, what we \textit{can} conclude is, given $\lbag Z_{0:t} \rbag = \lbag z_{0:t} \rbag$, that
\begin{align*}
    \lbag Z_{0:t} \rbag = \lbag z_{0:t} \rbag \implies \lbag L_{1:t-1} \rbag \in \Big\{\lbag l_{0:t\backslash i} \rbag \Big\}_{i\in \{0, ..., t\}},
\end{align*}
where $\lbag l_{0:t\backslash i} \rbag := \lbag l_{0:t} \rbag \backslash \{l_i\}$. Accordingly, conditioning on the event $\lbag Z_{0:t} \rbag = \lbag z_{0:t} \rbag$ implies the following equivalence:
\begin{align}\label{eq:cpc_oracle_empirical_inequality_equivalences}
    \beta^* \geq \hat{\beta}(\lbag L_{0:t-1}, B \rbag, \tilde{w}_{0:t-1}^{(\cdot)}, \tilde{w}_{\max}^{(\cdot)}, \alpha) \ \text{almost surely} \ & \iff \beta^* \geq \hat{\beta}(\lbag l_{0:t\backslash i}, B \rbag, \tilde{w}_{0:t-1}^{(\cdot)}, \tilde{w}_{\max}^{(\cdot)}, \alpha) \quad \forall \quad i \in \{0, ..., t\}.
\end{align}
So, we will focus on proving the latter. 




Recall that we have assumed that $l$ is bounded above by $B$, for any $\beta\in \mathbb{R}$, and by the definition of $\tilde{w}_{\max}^{(\cdot)}$, deterministically we have
\begin{align}\label{eq:cpc_step_3_conservative_adjustment}
    \sum_{j\in\{0, ..., t\}}l_j\cdot\tilde{w}_{j}^{(\beta)}
    \quad \leq \quad  B\cdot \tilde{w}_{\max}^{(\beta)} + \sum_{j\in\{0, ..., t\}\backslash i}l_j\cdot\tilde{w}_{j}^{(\beta)}  \qquad \forall \qquad i \in \{0, ..., t\},
\end{align}


where the right-hand side of the above inequalities can be interpreted as a conservative risk obtained by replacing $l_i$ with the bound $B$.


Thus, for any $\beta$, we can use \eqref{eq:cpc_step_3_conservative_adjustment} to conclude
\begin{align}\label{eq:cpc_step_3_risk_relation}
    B\cdot \tilde{w}_{\max}^{(\beta)} + \sum_{j\in\{0, ..., t\}\backslash i}l_j\cdot\tilde{w}_{i}^{(\beta)} \leq \alpha \implies & \sum_{j\in\{0, ..., t\}}l_j\cdot\tilde{w}_{j}^{(\beta)} \leq \alpha \qquad \forall \qquad i \in \{0, ..., t\}.
\end{align} 

So, by Eq. \eqref{eq:cpc_step_3_risk_relation} and the definition of $\hat{\beta}(\lbag l_{0:t\backslash i}, B \rbag)$, it must hold that 
\begin{align}\label{eq:cpc_pf_step_3_result}
    \beta^* \geq & \ \hat{\beta}(\lbag l_{0:t\backslash i}, B \rbag, \tilde{w}_{0:t-1}^{(\cdot)}, \tilde{w}_{\max}^{(\cdot)}, \alpha) \quad \forall \quad i \in \{0, ..., t\}.
\end{align}

\vspace{1cm}

\textbf{Step 4: Analyzing the bag-conditional risk of the empirical procedure via Lipschitz continuity}

Now we turn to plugging the empirical algorithm's hyperparameter into the bag-conditional risk, and then analyzing this risk using Lipschitz continuity and a type of stability called ``replace-one'' stability. This analysis is necessary because $\hat\beta=\hat{\beta}(\lbag L_{0:t-1}, B \rbag, \tilde{w}_{0:t-1}^{(\cdot)}, \tilde{w}_{\max}^{(\cdot)}, \alpha)$ is not constant conditional on $\lbag Z_{0:t} \rbag = \lbag z_{0:t} \rbag$, so Eq. \eqref{eq:cpc_step_2} would not directly apply. 

Begin by expanding the quantity $\mathbb{E}_{\pi_{0:t}^{(\hat\beta)}}\Big[L_{t} \mid \lbag Z_{0:t} \rbag = \lbag z_{0:t} \rbag\Big]$ by LOTE over the event $Z_{t}=z_i$:
\begin{align}\label{eq:cpc_empirical_bag_cond_LOTE}
    \mathbb{E}_{\pi_{0:t}^{(\hat\beta)}}\Big[&L_{t} \mid \lbag Z_{0:t} \rbag = \lbag z_{0:t} \rbag\Big] \nonumber \\
    & = \mathbb{E}_{\pi_{0:t}^{(\hat\beta)}}\Big[\mathbb{E}_{\pi_{0:t}(\hat{\beta}(\lbag L_{0:t-1} \rbag))}\Big[L_{t} \mid \lbag Z_{0:t} \rbag = \lbag z_{0:t} \rbag, Z_{t}=z_i\Big] \mid \lbag Z_{0:t} \rbag = \lbag z_{0:t} \rbag\Big] \nonumber \\
    & = \sum_{i=0}^{t}\mathbb{E}_{\pi_{0:t}^{(\hat\beta)}}\Big[L_{t} \mid \lbag Z_{0:t} \rbag = \lbag z_{0:t} \rbag, Z_{t}=z_i\Big]\cdot \mathbb{P}_{\pi_{0:t}^{(\hat\beta)}}\big(Z_{t}=z_i \mid \lbag Z_{0:t} \rbag = \lbag z_{0:t} \rbag\big) \nonumber \\
    & = \sum_{i=0}^{t}l_i\cdot \mathbb{P}_{\pi_{0:t}^{(\hat\beta)}}\big(Z_{t}=z_i \mid \lbag Z_{0:t} \rbag = \lbag z_{0:t} \rbag\big).
\end{align}


Writing out the conditional probability in Eq. \eqref{eq:cpc_empirical_bag_cond_LOTE} using the definition of conditional probability, and then (in both the numerator and denominator) the law of total probability and chain rule,
\begin{align*}
    & \mathbb{P}_{\pi_{0:t}^{(\hat\beta)}}\big(Z_{t}=z_i \mid \lbag Z_{0:t} \rbag = \lbag z_{0:t} \rbag\big) \\
    & = \frac{\sum_{\sigma:\sigma(t)=i}\pi_{0:t}^{(\hat\beta)}(z_{\sigma(0)}, ..., z_{\sigma(t)})}{\sum_{\sigma}\pi_{0:t}^{(\hat\beta)}(z_{\sigma(0)}, ..., z_{\sigma(t)})} \\
    & = \frac{\sum_{\sigma:\sigma(t)=i}\pi_{t\mid 0:t-1}^{(\hat\beta)}(z_{\sigma(t)} \mid z_{\sigma(0)}, ..., z_{\sigma(t-1)})\cdot \pi_{t-1\mid 0:t-2}(z_{\sigma(t-1)} \mid z_{\sigma(0)}, ..., z_{\sigma(t-2)})\cdots \pi_{0}(z_{\sigma(0)})}{\sum_{\sigma}\pi_{t\mid 0:t-1}^{(\hat\beta)}(z_{\sigma(t)} \mid z_{\sigma(0)}, ..., z_{\sigma(t-1)})\cdot \pi_{t-1\mid 0:t-2}(z_{\sigma(t-1)} \mid z_{\sigma(0)}, ..., z_{\sigma(t-2)})\cdots \pi_{0}(z_{\sigma(0)})},
\end{align*}
and by expanding the summation in the denominator, this is equivalent to 
\begin{align*}
    & \mathbb{P}_{\pi_{0:t}^{(\hat\beta)}}\big(Z_{t}=z_i \mid \lbag Z_{0:t} \rbag = \lbag z_{0:t} \rbag\big) \\
    & = \frac{\sum_{\sigma:\sigma(t)=i}\pi_{t\mid 0:t-1}^{(\hat\beta)}(z_{\sigma(t)} \mid z_{\sigma(0)}, ..., z_{\sigma(t-1)})\cdot \pi_{t-1\mid 0:t-2}(z_{\sigma(t-1)} \mid z_{\sigma(0)}, ..., z_{\sigma(t-2)})\cdots \pi_{0}(z_{\sigma(0)})}{\sum_{i}\sum_{\sigma:\sigma(t)=i}\pi_{t\mid 0:t-1}^{(\hat\beta)}(z_{\sigma(t)} \mid z_{\sigma(0)}, ..., z_{\sigma(t-1)})\cdot \pi_{t-1\mid 0:t-2}(z_{\sigma(t-1)} \mid z_{\sigma(0)}, ..., z_{\sigma(t-2)})\cdots \pi_{0}(z_{\sigma(0)})}.
\end{align*}
In this expanded form, examine the factor $\pi_{t\mid 0:t-1}^{(\hat\beta)}(z_{\sigma(t)} \mid z_{\sigma(0)}, ..., z_{\sigma(t-1)})$, and recalling that $\hat\beta:= \hat{\beta}(\lbag L_{0:t-1}, B \rbag, \tilde{w}_{0:t-1}^{(\cdot)}, \tilde{w}_{\max}^{(\cdot)}, \alpha)$, note that the conditioning on $z_{\sigma(0)}, ..., z_{\sigma(t-1)}$ inside the summation $\sum_{\sigma:\sigma(t)=i}$, where the permutations are $\sigma:\sigma(t)=i$, implies  $Z_t=z_{\sigma(t)}=z_i \implies \lbag L_{0:t-1} \rbag = \lbag l_{0:t\backslash i} \rbag$ (that is, implies a specific realization of $\lbag L_{0:t-1} \rbag$ among $t$ possible bags of observed losses). Accordingly, let us denote $\hat\beta^{\backslash i}:= \hat{\beta}(\lbag l_{0:t\backslash i}, B \rbag, \tilde{w}_{0:t-1}^{(\cdot)}, \tilde{w}_{\max}^{(\cdot)}, \alpha)$. Then, note that the numerator and the analogous inner summation in the denominator are the unnormalized weights in Eq. \eqref{eq:cp_weights_def_cpc_unnormed}, that is, we can concisely write those summations as $\sum_{\sigma:\sigma(t)=i}(\cdots) = w_i^{(\beta^{\backslash i})}$. That is, we have
\begin{align}
    & \mathbb{P}_{\pi_{0:t}^{(\hat\beta)}}\big(Z_{t}=z_i \mid \lbag Z_{0:t} \rbag = \lbag z_{0:t} \rbag\big)  = \frac{w_i^{(\beta^{\backslash i})}}{\sum_{j}w_j^{(\beta^{\backslash j})}}.
\end{align}

Plugging this into Eq. \eqref{eq:cpc_empirical_bag_cond_LOTE}, we have
\begin{align}\label{eq:empirical_risk_bag_conditional}
    \mathbb{E}_{\pi_{0:t}^{(\hat\beta)}}\Big[L_{t} \mid \lbag Z_{0:t} \rbag = \lbag z_{0:t} \rbag\Big] = 
    \sum_{i=0}^{t}l_i\cdot \frac{w_i^{(\beta^{\backslash i})}}{\sum_{j}w_j^{(\beta^{\backslash j})}} = \frac{1}{\sum_{j}w_j^{(\beta^{\backslash j})}}\sum_{i=0}^{t}l_i\cdot w_i^{(\beta^{\backslash i})}
\end{align}


Importantly, note that in Eq. \eqref{eq:empirical_risk_bag_conditional}, which can be viewed as the true bag-conditional risk of the empirical algorithm, that the
unnormalized weight in the numerator, $w_i^{(\beta^{\backslash i})}$, differs for each loss function, $l_i$, in the summation, and moreover the same is true in the denominator's summation, 
which is why we cannot apply Eq. \eqref{eq:cpc_step_2} to bound it directly. 


Instead, we now move toward relating Eq. \eqref{eq:empirical_risk_bag_conditional} and Eq. \eqref{eq:cpc_step_2} using Lipschitz continuity and the replace-one (in)stability of the algorithm. However, although in the analogous step assuming  parameterized losses (Section \ref{app:gcrc_proof_nonmonotonic_1dim}) we were able to directly apply Lipschitz continuity at this stage, here we must be more careful. That is, although in Section \ref{app:gcrc_proof_nonmonotonic_1dim} we could conservatively assume that each replace-one perturbation $\epsilon_i$ increased its corresponding loss by at most $K\epsilon_i$, a subtle difference here is that if we artificially increased
the unnormalized weight of certain (e.g., smaller-loss) indices by some amount, this may \textit{decrease} the overall bag-conditional weighted risk after re-normalizing the probability weights, and thus not result in a valid upper bound. In other words, conservative estimates of losses monotonically result in conservative estimates of the risk (in Section \ref{app:gcrc_proof_nonmonotonic_1dim}), but here, conservative estimates of (unnormalized) weights introduce an additional source of non-monotonicity that may not result in a conservative estimate of the weighted (bag-conditional) risk.






For all $i\in[n]$, denote the magnitude of perturbation to the hyperparameter $\hat{\beta}\big(\lbag l_{0:t\backslash i}, B \rbag, \tilde{w}_{0:t-1}^{(\cdot)}, \tilde{w}_{\max}^{(\cdot)}, \alpha\big)$ caused by replacing the loss function $l_i$ with $l_{t}$ as 
\begin{align*}
    \epsilon_i & = \Big|\hat{\beta}\big(\lbag l_{0:t-1}, B \rbag, \tilde{w}_{0:t-1}^{(\cdot)}, \tilde{w}_{\max}^{(\cdot)}, \alpha\big)-\hat{\beta}\big(\lbag l_{0:t\backslash i}, B \rbag, \tilde{w}_{0:t-1}^{(\cdot)}, \tilde{w}_{\max}^{(\cdot)}, \alpha\big)\Big| \\
    & = \Big|\hat{\beta}^{\backslash t}-\hat{\beta}^{\backslash i}\Big|.
\end{align*}

% By our assumption that the unnormalized weights times the loss, $l_i\cdot w_i^{(\beta)}$, are $K$-log-Lipschitz in $\beta$, note that
% \begin{align}\label{eq:K_log_Lipschitz_liwi}
%     \big|\log(l_i\cdot w_i^{(\hat\beta^{\backslash i})}) - \log(l_i\cdot w_i^{(\hat\beta^{\backslash t})})\big| = \bigg|\log\bigg(\frac{ w_i^{(\hat\beta^{\backslash i})}}{w_i^{(\hat\beta^{\backslash t})}}\bigg)\bigg| & \leq  K\epsilon_i \quad  \forall \quad i \in \{0, ..., t\} \nonumber \\
%     & \implies \frac{ w_i^{(\hat\beta^{\backslash i})}}{w_i^{(\hat\beta^{\backslash t})}} \leq e^{K\epsilon_i} \quad  \forall \quad i \in \{0, ..., t\}.
% \end{align}





% Now, although in the analogous step assuming  parameterized losses (Section \ref{app:gcrc_proof_nonmonotonic_1dim}) we were able to directly apply Lipschitz continuity, here we must be more careful. That is, although in Section \ref{app:gcrc_proof_nonmonotonic_1dim} we could conservatively assume that each replace-one perturbation $\epsilon_i$ increased its corresponding loss by at most $K\epsilon_i$, a subtle difference here is that if we artifically increased
% the unnormalized weight of certain (e.g., smaller-loss) indices by some amount, this may \textit{decrease} the overall bag-conditional weighted risk after re-normalizing the probability weights, and thus not result in a valid upper bound. In other words, conservative estimates of losses monotonically result in conservative estimates of the risk (in Section \ref{app:gcrc_proof_nonmonotonic_1dim}), but here, conservative estimates of (unnormalized) weights introduce an additional source of non-monotonicity that may not result in a conservative estimate of the weighted (bag-conditional) risk.
% Formally, 

% {
% \color{blue}
% \drew{re-check this equation after swapping assumptions}
% \begin{align*}
%     l_i\cdot w_i^{(\hat\beta^{\backslash i})} \leq l_i\cdot w_i^{(\hat\beta^{\backslash t})} + K\epsilon_i \quad & \forall \quad i \in \{0, ..., t\} \\
%     & \not\Rightarrow 
%     \qquad \frac{1}{\sum_{j}w_j^{(\beta^{\backslash j})}}\sum_{i=0}^{t}l_i\cdot w_i^{(\hat\beta^{\backslash i})} \leq \frac{1}{\sum_{j}(w_j^{(\beta^{\backslash t})} + K\epsilon_j/l_j)}\sum_{i=0}^{t}l_i w_i^{(\hat\beta^{\backslash t})}+ K\epsilon_i,
% \end{align*}
% where for emphasis, this can be roughly read as follows: conservative unnormalized weight estimates via Lipschitz continuity (LHS) do \textit{not} imply a valid risk upper bound (i.e., RHS would be desired but does not follow). 
% }

% For notational convenience, let $r_i:=e^{K\epsilon_i}$. 
% Then, note that by dividing this inequality by $l_i$, we have a similar bound for the unnormalized weights as
% \begin{align*}
%     w_i^{(\hat\beta^{\backslash i})} \leq w_i^{(\hat\beta^{\backslash t})} + Kr_i \quad & \forall \quad i \in \{0, ..., t\}.
% \end{align*}

% For notational convenience, let us denote $r_i:=e^{K\epsilon_i}$ for the upper bound on the ratio $w_i^{(\hat\beta^{\backslash i})}/w_i^{(\hat\beta^{\backslash t})}$ in Eq. \eqref{eq:K_log_Lipschitz_liwi}.



% We can still upper-bound the (bag-conditional) weighted risk, however, 
% by conservatively assuming that the largest loss values receive the largest weight increase. 
% To formalize this, without loss of generality hereon assume that the loss  are sorted so that their indices indicate their rank, that is, $i \leq j \iff l_i\leq l_j$ for all $i, j\in \{0, ..., t\}$ and $i\neq j$. 
% Separately, for notational convenience, let us denote $r_i:=w_i^{(\hat\beta^{\backslash i})}/w_i^{(\hat\beta^{\backslash t})}$, the ratio in how the unnormalized weight changes. 
% Plugging in $w_i^{(\hat\beta^{\backslash i})} = w_i^{(\hat\beta^{\backslash t})}r_i$ into Eq. \eqref{eq:empirical_risk_bag_conditional}, we equivalently have



% Note that we can rewrite $w_i^{(\hat\beta^{\backslash i})}$ as follows, by multiplying by $1 = \frac{w_i^{(\hat\beta^{\backslash t})}}{w_i^{(\hat\beta^{\backslash t})}}$ and then adding $0 = w_i^{(\hat\beta^{\backslash t})} - w_i^{(\hat\beta^{\backslash t})}$ in the numerator, to get
Making the substitution $w_i^{(\hat\beta^{\backslash i})} = w_i^{(\hat\beta^{\backslash i})} + (w_i^{(\hat\beta^{\backslash t})} - w_i^{(\hat\beta^{\backslash t})})$ in Eq. \eqref{eq:empirical_risk_bag_conditional} and then rearranging terms, we have
% \begin{align*}
%     w_i^{(\hat\beta^{\backslash i})} = w_i^{(\hat\beta^{\backslash t})}\Big(\frac{w_i^{(\hat\beta^{\backslash i})}}{w_i^{(\hat\beta^{\backslash t})}}\Big)= w_i^{(\hat\beta^{\backslash t})}\Big(\frac{w_i^{(\hat\beta^{\backslash i})} + w_i^{(\hat\beta^{\backslash t})} - w_i^{(\hat\beta^{\backslash t})}}{w_i^{(\hat\beta^{\backslash t})}}\Big) = w_i^{(\hat\beta^{\backslash t})}\Big(1 + \frac{w_i^{(\hat\beta^{\backslash i})} - w_i^{(\hat\beta^{\backslash t})}}{w_i^{(\hat\beta^{\backslash t})}}\Big).
% \end{align*}

% Substituting this into  Eq. \eqref{eq:empirical_risk_bag_conditional} and then rearranging terms, we equivalently have
\begin{align*}
% \label{eq:cpc_pf_centering_substitution}
    % \frac{1}{\sum_{j}w_j^{(\beta^{\backslash j})}}\sum_{i=0}^{t}l_i\cdot w_i^{(\beta^{\backslash i})} 
    \mathbb{E}_{\pi_{0:t}^{(\hat\beta)}}\Big[L_{t} \mid \lbag Z_{0:t} \rbag = \lbag z_{0:t} \rbag\Big] 
    % = & \frac{1}{\sum_{j}w_j^{(\hat\beta^{\backslash j})}}\sum_{i=0}^{t}l_i\cdot w_i^{(\hat\beta^{\backslash t})}\Big(1 + \frac{w_i^{(\hat\beta^{\backslash i})} - w_i^{(\hat\beta^{\backslash t})}}{w_i^{(\hat\beta^{\backslash t})}}\Big) \nonumber \\
    % = & \frac{1}{\sum_{j}w_j^{(\hat\beta^{\backslash j})}}\sum_{i=0}^{t}l_i\cdot w_i^{(\hat\beta^{\backslash t})} + l_i\cdot w_i^{(\hat\beta^{\backslash t})}\frac{w_i^{(\hat\beta^{\backslash i})} - w_i^{(\hat\beta^{\backslash t})}}{w_i^{(\hat\beta^{\backslash t})}} \nonumber \\
    = & \frac{1}{\sum_{j}w_j^{(\hat\beta^{\backslash j})}}\sum_{i=0}^{t}l_i\cdot w_i^{(\hat\beta^{\backslash t})} +  (l_i\cdot w_i^{(\hat\beta^{\backslash i})} - l_i\cdot w_i^{(\hat\beta^{\backslash t})}) \nonumber \\
    = & \frac{1}{\sum_{j}w_j^{(\hat\beta^{\backslash j})}}\sum_{i=0}^{t}l_i\cdot w_i^{(\hat\beta^{\backslash t})} +  \frac{1}{\sum_{j}w_j^{(\hat\beta^{\backslash j})}}\sum_{i=0}^{t}(l_i\cdot w_i^{(\hat\beta^{\backslash i})} - l_i\cdot w_i^{(\hat\beta^{\backslash t})}).
\end{align*}
% And then applying $K$-Lipschitz in $l_iw_i^{(\beta)}$,
% \begin{align*}
%     \frac{1}{\sum_{j}w_j^{(\beta^{\backslash j})}}\sum_{i=0}^{t}l_i\cdot w_i^{(\beta^{\backslash i})} \leq \frac{1}{\sum_{j}w_j^{(\hat\beta^{\backslash j})}}\sum_{i=0}^{t}l_i\cdot w_i^{(\hat\beta^{\backslash t})} +  K\epsilon_i.
% \end{align*}

Now multiplying the first term by $1 = \frac{\sum_{j}w_j^{(\hat\beta^{\backslash t})}}{\sum_{j}w_j^{(\hat\beta^{\backslash t})}}$,
\begin{align*}
    % \frac{1}{\sum_{j}w_j^{(\beta^{\backslash j})}}\sum_{i=0}^{t}l_i\cdot w_i^{(\beta^{\backslash i})} 
    \mathbb{E}_{\pi_{0:t}^{(\hat\beta)}}\Big[L_{t} \mid \lbag Z_{0:t} \rbag = \lbag z_{0:t} \rbag\Big] 
    = & \frac{\sum_{j}w_j^{(\hat\beta^{\backslash t})}}{\sum_{j}w_j^{(\hat\beta^{\backslash t})}}\frac{1}{\sum_{j}w_j^{(\hat\beta^{\backslash j})}}\sum_{i=0}^{t}l_i\cdot w_i^{(\hat\beta^{\backslash t})} +  \frac{1}{\sum_{j}w_j^{(\hat\beta^{\backslash j})}}\sum_{i=0}^{t}(l_i\cdot w_i^{(\hat\beta^{\backslash i})} - l_i\cdot w_i^{(\hat\beta^{\backslash t})}) \\
    = & \frac{\sum_{j}w_j^{(\hat\beta^{\backslash t})}}{\sum_{j}w_j^{(\hat\beta^{\backslash j})}}\sum_{i=0}^{t}\frac{l_i\cdot w_i^{(\hat\beta^{\backslash t})}}{\sum_{j}w_j^{(\hat\beta^{\backslash t})}} +  \frac{1}{\sum_{j}w_j^{(\hat\beta^{\backslash j})}}\sum_{i=0}^{t}(l_i\cdot w_i^{(\hat\beta^{\backslash i})} - l_i\cdot w_i^{(\hat\beta^{\backslash t})}),
\end{align*}

Then by Eq. \eqref{eq:cpc_step_2}, we know that $\sum_{i=0}^{t}\frac{l_i\cdot w_i^{(\hat\beta^{\backslash t})}}{\sum_{j}w_j^{(\hat\beta^{\backslash t})}} \leq \alpha$ (note that $\hat\beta^{\backslash t}$ is fixed with respect to the summations in the term), so
\begin{align}\label{eq:risk_bound_alpha_sub}
    % \frac{1}{\sum_{j}w_j^{(\beta^{\backslash j})}}\sum_{i=0}^{t}l_i\cdot w_i^{(\beta^{\backslash i})} 
    \mathbb{E}_{\pi_{0:t}^{(\hat\beta)}}\Big[L_{t} \mid \lbag Z_{0:t} \rbag = \lbag z_{0:t} \rbag\Big] 
    \leq & \ \frac{\sum_{j}w_j^{(\hat\beta^{\backslash t})}}{\sum_{j}w_j^{(\hat\beta^{\backslash j})}}\cdot\alpha +  \frac{1}{\sum_{j}w_j^{(\hat\beta^{\backslash j})}}\sum_{i=0}^{t}(l_i\cdot w_i^{(\hat\beta^{\backslash i})} - l_i\cdot w_i^{(\hat\beta^{\backslash t})}).
\end{align}

Plugging in $w_j^{(\hat\beta^{\backslash t})} = w_j^{(\hat\beta^{\backslash t})} + (w_j^{(\hat\beta^{\backslash j})} - w_j^{(\hat\beta^{\backslash j})})$ into the numerator of the first term of the RHS of \eqref{eq:risk_bound_alpha_sub} and rearranging terms,

% \begin{align*}
%     \text{RHS} 
%     = & \frac{\sum_{j}(w_j^{(\hat\beta^{\backslash t})} + w_j^{(\hat\beta^{\backslash j})} - w_j^{(\hat\beta^{\backslash j})})}{\sum_{j}w_j^{(\hat\beta^{\backslash j})}}\cdot\alpha +  \frac{1}{\sum_{j}w_j^{(\hat\beta^{\backslash j})}}\sum_{i=0}^{t}(l_i\cdot w_i^{(\hat\beta^{\backslash i})} - l_i\cdot w_i^{(\hat\beta^{\backslash t})}) \\
%     = & \frac{\sum_{j}w_j^{(\hat\beta^{\backslash j})} + \sum_{j}(w_j^{(\hat\beta^{\backslash t})} - w_j^{(\hat\beta^{\backslash j})})}{\sum_{j}w_j^{(\hat\beta^{\backslash j})}}\cdot\alpha +  \frac{1}{\sum_{j}w_j^{(\hat\beta^{\backslash j})}}\sum_{i=0}^{t}(l_i\cdot w_i^{(\hat\beta^{\backslash i})} - l_i\cdot w_i^{(\hat\beta^{\backslash t})}) \\
%     = & \bigg(1 + \frac{\sum_{j}(w_j^{(\hat\beta^{\backslash t})} - w_j^{(\hat\beta^{\backslash j})})}{\sum_{j}w_j^{(\hat\beta^{\backslash j})}}\bigg)\cdot\alpha +  \frac{1}{\sum_{j}w_j^{(\hat\beta^{\backslash j})}}\sum_{i=0}^{t}(l_i\cdot w_i^{(\hat\beta^{\backslash i})} - l_i\cdot w_i^{(\hat\beta^{\backslash t})}) \\ 
%     = & \ \alpha + \frac{\sum_{j}(w_j^{(\hat\beta^{\backslash t})} - w_j^{(\hat\beta^{\backslash j})})}{\sum_{j}w_j^{(\hat\beta^{\backslash j})}}\cdot\alpha +  \frac{1}{\sum_{j}w_j^{(\hat\beta^{\backslash j})}}\sum_{i=0}^{t}(l_i\cdot w_i^{(\hat\beta^{\backslash i})} - l_i\cdot w_i^{(\hat\beta^{\backslash t})}) \\
%     = & \ \alpha + \frac{\sum_{i=0}^{t}(\alpha\cdot w_i^{(\hat\beta^{\backslash t})} - \alpha\cdot w_i^{(\hat\beta^{\backslash i})}) + \sum_{i=0}^{t}(l_i\cdot w_i^{(\hat\beta^{\backslash i})} - l_i\cdot w_i^{(\hat\beta^{\backslash t})})}{\sum_{j}w_j^{(\hat\beta^{\backslash j})}} \\
%     = & \ \alpha + \frac{\sum_{i=0}^{t}\alpha\cdot w_i^{(\hat\beta^{\backslash t})} - \alpha\cdot w_i^{(\hat\beta^{\backslash i})} + l_i\cdot w_i^{(\hat\beta^{\backslash i})} - l_i\cdot w_i^{(\hat\beta^{\backslash t})}}{\sum_{j}w_j^{(\hat\beta^{\backslash j})}} \\
%     = & \ \alpha + \frac{\sum_{i=0}^{t}l_i\cdot w_i^{(\hat\beta^{\backslash i})} - \alpha\cdot w_i^{(\hat\beta^{\backslash i})} + \alpha\cdot w_i^{(\hat\beta^{\backslash t})} - l_i\cdot w_i^{(\hat\beta^{\backslash t})}}{\sum_{j}w_j^{(\hat\beta^{\backslash j})}} \\
%     = & \ \alpha + \frac{\sum_{i=0}^{t}w_i^{(\hat\beta^{\backslash i})}(l_i  - \alpha) - w_i^{(\hat\beta^{\backslash t})}(l_i -\alpha)}{\sum_{j}w_j^{(\hat\beta^{\backslash j})}} \\
%     = & \ \alpha + \frac{\sum_{i=0}^{t}(l_i -\alpha)(w_i^{(\hat\beta^{\backslash i})} - w_i^{(\hat\beta^{\backslash t})})}{\sum_{j}w_j^{(\hat\beta^{\backslash j})}}
% \end{align*}

% Partition the indices based on the sign of $l_i -\alpha$ (i.e., based on whether $l_i \geq \alpha$):
% \begin{align*}
%     \mathcal{I}_+ & := \big\{i\in \{0, ..., t\} : l_i \geq \alpha\big\} \\
%     \mathcal{I}_- & := \big\{i\in \{0, ..., t\} : l_i < \alpha \big\},
% \end{align*}

% then, 
% \begin{align*}
%     \text{RHS} & = \ \alpha + \frac{\sum_{i\in\mathcal{I}_+}(l_i -\alpha)(w_i^{(\hat\beta^{\backslash i})} - w_i^{(\hat\beta^{\backslash t})}) + \sum_{i\in\mathcal{I}_-}(l_i -\alpha)(w_i^{(\hat\beta^{\backslash i})} - w_i^{(\hat\beta^{\backslash t})})}{\sum_{j}w_j^{(\hat\beta^{\backslash j})}} \\
%     & = \ \alpha + \frac{\sum_{i\in\mathcal{I}_+}|l_i -\alpha|(w_i^{(\hat\beta^{\backslash i})} - w_i^{(\hat\beta^{\backslash t})}) - \sum_{i\in\mathcal{I}_-}|l_i -\alpha|(w_i^{(\hat\beta^{\backslash i})} - w_i^{(\hat\beta^{\backslash t})})}{\sum_{j}w_j^{(\hat\beta^{\backslash j})}}
% \end{align*}

% Notice then that the first sum, $\sum_{i\in\mathcal{I}_+}|l_i -\alpha|(w_i^{(\hat\beta^{\backslash i})} - w_i^{(\hat\beta^{\backslash t})})$, is monotonically increasing in 

% \begin{align*}
%     \frac{\sum_{j}w_j^{(\hat\beta^{\backslash t})}}{\sum_{j}w_j^{(\hat\beta^{\backslash j})}}\cdot\alpha & +  \frac{1}{\sum_{j}w_j^{(\hat\beta^{\backslash j})}}\sum_{i=0}^{t}(l_i\cdot w_i^{(\hat\beta^{\backslash i})} - l_i\cdot w_i^{(\hat\beta^{\backslash t})}) \\
%     = & \frac{\sum_{j}(w_j^{(\hat\beta^{\backslash t})} + w_j^{(\hat\beta^{\backslash j})} - w_j^{(\hat\beta^{\backslash j})})}{\sum_{j}w_j^{(\hat\beta^{\backslash j})}}\cdot\alpha +  \frac{1}{\sum_{j}w_j^{(\hat\beta^{\backslash j})}}\sum_{i=0}^{t}(l_i\cdot w_i^{(\hat\beta^{\backslash i})} - l_i\cdot w_i^{(\hat\beta^{\backslash t})}) \\
%     = & \frac{\sum_{j}w_j^{(\hat\beta^{\backslash j})} + \sum_{j}(w_j^{(\hat\beta^{\backslash t})} - w_j^{(\hat\beta^{\backslash j})})}{\sum_{j}w_j^{(\hat\beta^{\backslash j})}}\cdot\alpha +  \frac{1}{\sum_{j}w_j^{(\hat\beta^{\backslash j})}}\sum_{i=0}^{t}(l_i\cdot w_i^{(\hat\beta^{\backslash i})} - l_i\cdot w_i^{(\hat\beta^{\backslash t})}) \\
%     = & \bigg(1 + \frac{\sum_{j}(w_j^{(\hat\beta^{\backslash t})} - w_j^{(\hat\beta^{\backslash j})})}{\sum_{j}w_j^{(\hat\beta^{\backslash j})}}\bigg)\cdot\alpha +  \frac{1}{\sum_{j}w_j^{(\hat\beta^{\backslash j})}}\sum_{i=0}^{t}(l_i\cdot w_i^{(\hat\beta^{\backslash i})} - l_i\cdot w_i^{(\hat\beta^{\backslash t})}) \\ 
%     = & \ \alpha + \frac{\sum_{j}(w_j^{(\hat\beta^{\backslash t})} - w_j^{(\hat\beta^{\backslash j})})}{\sum_{j}w_j^{(\hat\beta^{\backslash j})}}\cdot\alpha +  \frac{1}{\sum_{j}w_j^{(\hat\beta^{\backslash j})}}\sum_{i=0}^{t}(l_i\cdot w_i^{(\hat\beta^{\backslash i})} - l_i\cdot w_i^{(\hat\beta^{\backslash t})}) \\
%     = & \ \alpha + \frac{\sum_{i=0}^{t}(\alpha\cdot w_i^{(\hat\beta^{\backslash t})} - \alpha\cdot w_i^{(\hat\beta^{\backslash i})}) + \sum_{i=0}^{t}(l_i\cdot w_i^{(\hat\beta^{\backslash i})} - l_i\cdot w_i^{(\hat\beta^{\backslash t})})}{\sum_{j}w_j^{(\hat\beta^{\backslash j})}} \\
%     = & \ \alpha + \frac{\sum_{i=0}^{t}\alpha\cdot w_i^{(\hat\beta^{\backslash t})} - \alpha\cdot w_i^{(\hat\beta^{\backslash i})} + l_i\cdot w_i^{(\hat\beta^{\backslash i})} - l_i\cdot w_i^{(\hat\beta^{\backslash t})}}{\sum_{j}w_j^{(\hat\beta^{\backslash j})}} \\
%     = & \ \alpha + \frac{\sum_{i=0}^{t}\alpha\cdot (w_i^{(\hat\beta^{\backslash t})} -  w_i^{(\hat\beta^{\backslash i})}) + l_i\cdot (w_i^{(\hat\beta^{\backslash i})} - w_i^{(\hat\beta^{\backslash t})})}{\sum_{j}w_j^{(\hat\beta^{\backslash j})}} 
%     % = & \ \alpha + \frac{\sum_{i=0}^{t}\alpha\cdot (w_i^{(\hat\beta^{\backslash t})} -  w_i^{(\hat\beta^{\backslash i})}) + l_i\cdot (w_i^{(\hat\beta^{\backslash i})} - w_i^{(\hat\beta^{\backslash t})})}{\sum_{j}w_j^{(\hat\beta^{\backslash j})}}
%     % = & \ \alpha + \frac{\sum_{i=0}^{t}l_i\cdot w_i^{(\hat\beta^{\backslash i})} - \alpha\cdot w_i^{(\hat\beta^{\backslash i})} + \alpha\cdot w_i^{(\hat\beta^{\backslash t})} - l_i\cdot w_i^{(\hat\beta^{\backslash t})}}{\sum_{j}w_j^{(\hat\beta^{\backslash j})}} \\
%     % = & \ \alpha + \frac{\sum_{i=0}^{t}w_i^{(\hat\beta^{\backslash i})}(l_i  - \alpha) - w_i^{(\hat\beta^{\backslash t})}(l_i -\alpha)}{\sum_{j}w_j^{(\hat\beta^{\backslash j})}} \\
%     % = & \ \alpha + \frac{\sum_{i=0}^{t}(l_i -\alpha)(w_i^{(\hat\beta^{\backslash i})} - w_i^{(\hat\beta^{\backslash t})})}{\sum_{j}w_j^{(\hat\beta^{\backslash j})}}
% \end{align*}

% Assuming that $w_i^{(\beta)}$ is $K_w$-Lipschitz and that $l_iw_i^{(\beta)}$ is $K_l$-Lipschitz, then,

% \begin{align*}
%     \frac{\sum_{j}w_j^{(\hat\beta^{\backslash t})}}{\sum_{j}w_j^{(\hat\beta^{\backslash j})}}\cdot\alpha  +  \frac{1}{\sum_{j}w_j^{(\hat\beta^{\backslash j})}}\sum_{i=0}^{t}(l_i\cdot w_i^{(\hat\beta^{\backslash i})} - l_i\cdot w_i^{(\hat\beta^{\backslash t})}) \leq  & \ \alpha + \frac{\sum_{i=0}^{t}\alpha\cdot K_w\epsilon_i + K_l\epsilon_i}{\sum_{j}w_j^{(\hat\beta^{\backslash j})}} \\
%     & = \ \alpha + \frac{\alpha\cdot K_w\sum_{i=0}^{t}\epsilon_i + K_{l}\sum_{i=0}^{t}\epsilon_i}{\sum_{j}w_j^{(\hat\beta^{\backslash j})}} 
% \end{align*}

% and after marginalizing, 

% \begin{align*}
%     \mathbb{E}[L_t] 
%     % & \leq \alpha +  \alpha K_w\epsilon + K_l\epsilon \\
%     & \leq  \alpha(1 +  K_w\epsilon) + K_l\epsilon
% \end{align*}

% Note that
% \begin{align*}
%     w_i^{(\hat\beta^{\backslash t})} = w_i^{(\hat\beta^{\backslash i})}\Big(\frac{w_i^{(\hat\beta^{\backslash t})}}{w_i^{(\hat\beta^{\backslash i})}}\Big)= w_i^{(\hat\beta^{\backslash i})}\Big(\frac{w_i^{(\hat\beta^{\backslash t})} + w_i^{(\hat\beta^{\backslash i})} - w_i^{(\hat\beta^{\backslash i})}}{w_i^{(\hat\beta^{\backslash i})}}\Big) = w_i^{(\hat\beta^{\backslash i})}\Big(1 + \frac{w_i^{(\hat\beta^{\backslash t})} - w_i^{(\hat\beta^{\backslash i})}}{w_i^{(\hat\beta^{\backslash i})}}\Big).
% \end{align*}

% \begin{align*}
%     \text{RHS} = & \frac{\sum_{j}w_j^{(\hat\beta^{\backslash t})}}{\sum_{j}w_j^{(\hat\beta^{\backslash t})}}\frac{1}{\sum_{j}w_j^{(\hat\beta^{\backslash j})}}\sum_{i=0}^{t}l_i\cdot w_i^{(\hat\beta^{\backslash t})} +  (l_i\cdot w_i^{(\hat\beta^{\backslash i})} - l_i\cdot w_i^{(\hat\beta^{\backslash t})})  \\
%     = & \frac{\sum_{j}w_j^{(\hat\beta^{\backslash t})}}{\sum_{j}w_j^{(\hat\beta^{\backslash t})}}\frac{1}{\sum_{j}w_j^{(\hat\beta^{\backslash j})}}\bigg(\sum_{i=0}^{t}l_i\cdot w_i^{(\hat\beta^{\backslash t})} +  \sum_{i=0}^{t}(l_i\cdot w_i^{(\hat\beta^{\backslash i})} - l_i\cdot w_i^{(\hat\beta^{\backslash t})}) \bigg) \\
%     = & \frac{\sum_{j}w_j^{(\hat\beta^{\backslash t})}}{\sum_{j}w_j^{(\hat\beta^{\backslash t})}}\frac{1}{\sum_{j}w_j^{(\hat\beta^{\backslash j})}}\sum_{i=0}^{t}l_i\cdot w_i^{(\hat\beta^{\backslash t})} +  \frac{\sum_{j}w_j^{(\hat\beta^{\backslash t})}}{\sum_{j}w_j^{(\hat\beta^{\backslash t})}}\frac{1}{\sum_{j}w_j^{(\hat\beta^{\backslash j})}}\sum_{i=0}^{t}(l_i\cdot w_i^{(\hat\beta^{\backslash i})} - l_i\cdot w_i^{(\hat\beta^{\backslash t})}) 
% \end{align*}

% Now making the substitution $w_j^{(\hat\beta^{\backslash j})} = w_j^{(\hat\beta^{\backslash j})} + w_j^{(\hat\beta^{\backslash t})} - w_j^{(\hat\beta^{\backslash t})}$
% in the denominator on the RHS,
% \begin{align*}
%     \text{RHS} = & \frac{1}{\sum_{j}w_j^{(\hat\beta^{\backslash j})} + w_j^{(\hat\beta^{\backslash t})} - w_j^{(\hat\beta^{\backslash t})}}\sum_{i=0}^{t}l_i\cdot w_i^{(\hat\beta^{\backslash t})} +  K\epsilon_i \\
%     = & \frac{1}{\sum_{j}w_j^{(\hat\beta^{\backslash t})}\Big(1 + \frac{w_i^{(\hat\beta^{\backslash j})} - w_j^{(\hat\beta^{\backslash t})}}{w_j^{(\hat\beta^{\backslash t})}}\Big)}\sum_{i=0}^{t}l_i\cdot w_i^{(\hat\beta^{\backslash t})} +  K\epsilon_i 
% \end{align*}


% \begin{align*}
%     w_i^{(\hat\beta^{\backslash t})}\Big(1 + \frac{w_i^{(\hat\beta^{\backslash i})} - w_i^{(\hat\beta^{\backslash t})}}{w_i^{(\hat\beta^{\backslash t})}}\Big) & = w_i^{(\hat\beta^{\backslash t})}\Big(\frac{w_i^{(\hat\beta^{\backslash t})}}{w_i^{(\hat\beta^{\backslash t})}} + \frac{w_i^{(\hat\beta^{\backslash i})} - w_i^{(\hat\beta^{\backslash t})}}{w_i^{(\hat\beta^{\backslash t})}}\Big) \\
%     & = w_i^{(\hat\beta^{\backslash t})}\Big(\frac{w_i^{(\hat\beta^{\backslash i})}}{w_i^{(\hat\beta^{\backslash t})}}\Big)
% \end{align*}
% \begin{align*}
%     r_i = \frac{w_i^{(\hat\beta^{\backslash i})}}{w_i^{(\hat\beta^{\backslash t})}} = \frac{w_i^{(\hat\beta^{\backslash i})} + w_i^{(\hat\beta^{\backslash t})} - w_i^{(\hat\beta^{\backslash t})}}{w_i^{(\hat\beta^{\backslash t})}} = 1 + \frac{w_i^{(\hat\beta^{\backslash i})} - w_i^{(\hat\beta^{\backslash t})}}{w_i^{(\hat\beta^{\backslash t})}}
% \end{align*}

% Now, let $\sigma'$ denote a permutation of the indices that sorts the $\{r_i\}_{i=0}^t$ in non-decreasing order, that is, 
% $\sigma'(i) \leq \sigma'(j) \iff r_{\sigma'(i)} \leq r_{\sigma'(j)}$ for all $i, j\in \{0, ..., t\}$ and $i\neq j$. 



% A given point's change in normalized weight is
% \begin{align*}
%     \frac{w_i^{(\hat\beta^{\backslash t})}r_{\sigma'(i)}}{\sum_k w_k^{(\hat\beta^{\backslash t})}r_{\sigma'(k)}} - \frac{w_i^{(\hat\beta^{\backslash t})}}{\sum_kw_k^{(\hat\beta^{\backslash t})}} = w_i^{(\hat\beta^{\backslash t})}\bigg(\frac{r_{\sigma'(i)}}{\sum_kw_k^{(\hat\beta^{\backslash t})}r_{\sigma'(k)}} - \frac{1}{\sum_k w_k^{(\hat\beta^{\backslash t})}}\bigg)
% \end{align*}

% Comparing this quantity for two 

% % \begin{align*}
% %     \quad \implies \quad \frac{w_i^{(\hat\beta^{\backslash i})}}{\sum_j w_j^{(\hat\beta^{\backslash j})}} = \frac{w_i^{(\hat\beta^{\backslash t})}r_i}{\sum_j w_j^{(\hat\beta^{\backslash t})}r_j} \leq 
% % \end{align*}
% % Now, note that the permutation $\sigma'$ also sorts the re-normalized weights, that is 

% % Separately, let $\sigma'$ denote a permutation of the indices that sorts the perturbation magnitudes, $\{\epsilon_i\}_{i=0}^t$, in non-decreasing order, that is, 
% % $\sigma'(i) \leq \sigma'(j) \iff \epsilon_{\sigma'(i)} \leq \epsilon_{\sigma'(j)}$ for all $i, j\in \{0, ..., t\}$ and $i\neq j$. 

% So,
% \begin{align*}
%     l_i\cdot w_i^{(\hat\beta^{\backslash i})} \leq l_{i}\cdot w_i^{(\hat\beta^{\backslash t})}\cdot r_{\sigma'(i)} \quad \forall \quad i \in \{0, ..., t\}
% \end{align*}
% and
% \begin{align*}
%     \sum_{j}l_j\cdot w_j^{(\beta^{\backslash j})} \leq \sum_{j}l_j\cdot w_j^{(\beta^{\backslash t})}\cdot r_{\sigma'(j)}
% \end{align*}
% % \begin{align*}
% %     \frac{1}{\sum_{j}w_j^{(\beta^{\backslash j})}}l_t\cdot w_t^{(\hat\beta^{\backslash i})} \leq \frac{1}{\sum_{j}(w_j^{(\beta^{\backslash t})}\cdot r_{\sigma'(j)})}l_{t}\cdot w_t^{(\hat\beta^{\backslash t})}\cdot r_{\sigma'(t)}
% % \end{align*}

% % \begin{align*}
% %     \epsilon_{\sigma'(i)} \leq &\epsilon_{\sigma'(j)} \quad \forall \quad i, j\in \{0, ..., t\}, i\neq j \qquad \nonumber \\
% %     &\implies \qquad \frac{\sum_{j}w_j^{(\beta^{\backslash j})}}{\sum_{j}(w_j^{(\beta^{\backslash t})} + K\epsilon_{\sigma'(j)}/l_j)}K\epsilon_{\sigma'(i)} \leq \frac{\sum_{j}w_j^{(\beta^{\backslash j})}}{\sum_{j}(w_j^{(\beta^{\backslash t})} + K\epsilon_{\sigma'(j)}/l_j)}K\epsilon_{\sigma'(j)} \quad \forall \quad i, j\in \{0, ..., t\}, i\neq j
% % \end{align*}


% Then, we \textit{can} conclude 
% % For notation to formalize this, let $\sigma_l$ denote a permutation of the indices that sorts the losses, $\{l_i\}_{i=0}^t$, in ascending order, that is, 
% % $\sigma_l(i) \leq \sigma_l(j) \iff l_{\sigma_l(i)} \leq l_{\sigma_l(j)}$ for all $i, j\in \{0, ..., t\}$ and $i\neq j$. Similarly, let $\sigma'$ denote a permutation of the indices that sorts the perturbation magnitudes, $\{\epsilon_i\}_{i=0}^t$, in ascending order, that is, 
% % $\sigma'(i) \leq \sigma'(j) \iff \epsilon_{\sigma'(i)} \leq \epsilon_{\sigma'(j)}$ for all $i, j\in \{0, ..., t\}$ and $i\neq j$. Then, we \textit{can} conclude 
% \begin{align}
%     l_i\cdot w_i^{(\hat\beta^{\backslash i})} \leq l_i\cdot w_i^{(\hat\beta^{\backslash t})} + K\epsilon_{\sigma'(i)} & \quad \forall \quad i \in \{0, ..., t\} \nonumber \\
%     & \implies \quad \frac{1}{\sum_{j}w_j^{(\beta^{\backslash j})}}\sum_{i=0}^{t}l_i\cdot w_i^{(\hat\beta^{\backslash i})} \leq \frac{1}{\sum_{j}(w_j^{(\beta^{\backslash t})} + K\epsilon_{\sigma'(j)}/l_j)}\sum_{i=0}^{t}l_{i}w_i^{(\hat\beta^{\backslash t})}+ K\epsilon_{\sigma'(i)},
% \end{align}
% % using $K$-Lipschitz continuity of the $w_i^{(\beta)}$ in $\beta$ 
% where the LHS of the implication is by Lipschitz continuity and the RHS of the implication can be interpreted as bounding the bag-conditional weighted empirical risk above by the conservative  estimate one would obtain by moving probability mass increments (obtained from Lipschitz continuity) toward the largest loss values (i.e., increasing the unnormalized weight of the largest loss value, $l_t$, by the largest resulting increment, $K\epsilon_{\sigma'(t)}$, and so on). 

% Now, consider the normalization term in the upper bound, $1/(\sum_{j}w_j^{(\beta^{\backslash t})}+K\epsilon_{\sigma'(j)}/l_j)$. Note that $K\epsilon_{\sigma'(j)}/l_j$ is nonnegative (we assume $l_j>0$), so $1/(\sum_{j}w_j^{(\beta^{\backslash t})}+K\epsilon_{\sigma'(j)}/l_j) \leq 1/\sum_{j}w_j^{(\beta^{\backslash t})}$ and we have
% % Expanding the RHS of the implication,

% \begin{align*}
%     \frac{1}{\sum_{j}w_j^{(\beta^{\backslash j})}}\sum_{i=0}^{t}l_i\cdot w_i^{(\hat\beta^{\backslash i})} \leq &  \frac{1}{\sum_{j}w_j^{(\beta^{\backslash t})}}\Big(\sum_{i=0}^{t}l_{t}w_i^{(\hat\beta^{\backslash t})}+ K\epsilon_{\sigma'(i)}\Big) \\
%     \leq &  \frac{1}{\sum_{j}w_j^{(\beta^{\backslash t})}}\sum_{i=0}^{t}l_{t}w_i^{(\hat\beta^{\backslash t})}+ \frac{1}{\sum_{j}w_j^{(\beta^{\backslash t})}}\sum_{i=0}^{t}K\epsilon_{\sigma'(i)}
% \end{align*}





% \begin{align*}
%     \frac{1}{\sum_{j}w_j^{(\beta^{\backslash j})}}\sum_{i=0}^{t}l_i\cdot w_i^{(\hat\beta^{\backslash i})}  \leq \frac{1}{\sum_{j}w_j^{(\beta^{\backslash j})}}\sum_{i=0}^{t}l_i\cdot (w_i^{(\hat\beta^{\backslash t})} + K \epsilon_i)  = & \frac{1}{\sum_{j}w_j^{(\beta^{\backslash j})}}\sum_{i=0}^{t}l_i\cdot w_i^{(\hat\beta^{\backslash t})} + \frac{1}{\sum_{j}w_j^{(\beta^{\backslash j})}}\sum_{i=0}^{t}l_iK \epsilon_i \\
%     % \leq & \sum_{i=0}^{t}l_i\cdot \tilde{w}_i^{(\hat\beta^{\backslash t})} + \sum_{i=0}^{t}BK \epsilon_i
% \end{align*}

% Want to bound TV distance of PMF defined by normalized weights relative to replace-one perturbation:
% \begin{align*}
%     \frac{1}{2}\sum_{i}\Big|\frac{w_i^{(\beta^{\backslash i})}}{\sum_jw_j^{(\beta^{\backslash j})}} - \frac{w_i^{(\beta^{\backslash i})}}{\sum_j(w_j^{(\beta^{\backslash j})})}+K\epsilon_j\Big|
% \end{align*}
% \begin{align*}
%     \frac{1}{2}\sum_{i}\Big|\frac{w_i^{(\beta^{\backslash i})}}{\sum_jw_j^{(\beta^{\backslash j})}} - \frac{w_i^{(\beta^{\backslash i})}}{\sum_j(w_j^{(\beta^{\backslash j})})}+K\epsilon_j\Big|
% \end{align*}


% Partition indices based on whether their loss values exceed the weighted empirical mean, $\sum_{i=0}^{t}\frac{l_i\cdot w_i^{(\beta^{\backslash i})}}{\sum_{j}w_j^{(\beta^{\backslash j})}}$:
% \begin{align*}
%     \mathcal{I}_l := \big\{i\in \{0, ..., t\} : l_i < \sum_{i=0}^{t}\frac{l_i\cdot w_i^{(\beta^{\backslash i})}}{\sum_{j}w_j^{(\beta^{\backslash j})}} \big\} \\
%     \mathcal{I}_u := \big\{i\in \{0, ..., t\} : l_i \geq \sum_{i=0}^{t}\frac{l_i\cdot w_i^{(\beta^{\backslash i})}}{\sum_{j}w_j^{(\beta^{\backslash j})}}\big\}.
% \end{align*}
% For intuition, note that moving normalized weight or probability mass from $\mathcal{I}_l$ to $\mathcal{I}_u$
% increases the mean. 

% Separating the sums in both the numerator and denominator based on this partition,
% \begin{align*}
%     \frac{\sum_{i=0}^{t}l_i\cdot w_i^{(\beta^{\backslash i})}}{\sum_{j}w_j^{(\beta^{\backslash j})}} = & \frac{\Big(\sum_{i\in \mathcal{I}_l}l_i\cdot w_i^{(\beta^{\backslash i})} + \sum_{i\in \mathcal{I}_u}l_i\cdot w_i^{(\beta^{\backslash i})}\Big)}{\sum_{j\in \mathcal{I}_l}w_j^{(\beta^{\backslash j})} + \sum_{j\in\mathcal{I}_u}w_j^{(\beta^{\backslash j})}} 
%     % = & \sum_{i\in \mathcal{I}_l}\frac{l_i\cdot w_i^{(\beta^{\backslash i})}}{\sum_{j}w_j^{(\beta^{\backslash j})}} + \sum_{i\in \mathcal{I}_u}\frac{l_i\cdot w_i^{(\beta^{\backslash i})}}{\sum_{j}w_j^{(\beta^{\backslash j})}},
% \end{align*}

% and applying $K$-Lipschitz in $w_i^{(\beta^{\backslash i})}$,
% \begin{align*}
%     \frac{1}{\sum_{j}w_j^{(\beta^{\backslash j})}}\sum_{i=0}^{t}l_i\cdot w_i^{(\beta^{\backslash i})} \leq & \frac{\Big(\sum_{i\in \mathcal{I}_l}l_i\cdot (w_i^{(\beta^{\backslash i})}-K\epsilon_i) + \sum_{i\in \mathcal{I}_u}l_i\cdot (w_i^{(\beta^{\backslash i})}+K\epsilon_i)\Big)}{\sum_{j\in \mathcal{I}_l}(w_j^{(\beta^{\backslash j})}- K\epsilon_j) + \sum_{j\in\mathcal{I}_u}(w_j^{(\beta^{\backslash j})}+ K\epsilon_j)} 
% \end{align*}
% \begin{align*}
%     \sum_{i=0}^{t}\frac{l_i\cdot w_i^{(\beta^{\backslash i})}}{\sum_{j}w_j^{(\beta^{\backslash j})}} = & \sum_{i\in \mathcal{I}_l}\frac{l_i\cdot w_i^{(\beta^{\backslash i})}}{\sum_{j}w_j^{(\beta^{\backslash j})}} + \sum_{i\in \mathcal{I}_u}\frac{l_i\cdot w_i^{(\beta^{\backslash i})}}{\sum_{j}w_j^{(\beta^{\backslash j})}} \\
%     = & \sum_{i\in \mathcal{I}_l}\frac{l_i\cdot w_i^{(\beta^{\backslash i})}}{\sum_{j}w_j^{(\beta^{\backslash j})}} + \sum_{i\in \mathcal{I}_u}\frac{l_i\cdot w_i^{(\beta^{\backslash i})}}{\sum_{j}w_j^{(\beta^{\backslash j})}},
% \end{align*}


% {\color{red}
% Assuming that $l_i\cdot w_i$ is Lipschitz:

% By $K$-Lipschitz continuity, we know that for all $i\in \{0, ..., t\}$, that
% \begin{align*}
%     \Big|l_i\cdot w_i^{(\hat\beta^{\backslash t})} -l_i\cdot w_i^{(\hat\beta^{\backslash i})}\Big| \leq K \epsilon_i.
% \end{align*}


% So, we have
% % \begin{align*}
% %     \frac{1}{\sum_{i=0}^tw_i^{(\hat\beta^{\backslash i})}}\sum_{i=0}^{t}l_i\cdot w_i^{(\hat\beta^{\backslash i})} & \leq \frac{1}{\sum_{i=0}^tw_i^{(\hat\beta^{\backslash i})}}\sum_{i=0}^{t}(l_i\cdot w_i^{(\hat\beta^{\backslash t})} + K \epsilon_i)  = \frac{1}{\sum_{i=0}^tw_i^{(\hat\beta^{\backslash i})}}\sum_{i=0}^{t}l_i\cdot w_i^{(\hat\beta^{\backslash t})} + \frac{1}{\sum_{i=0}^tw_i^{(\hat\beta^{\backslash i})}}\sum_{i=0}^{t}K \epsilon_i
% % \end{align*}
% \begin{align*}
%     \frac{1}{\sum_{j}w_j^{(\beta^{\backslash j})}}\sum_{i=0}^{t}l_i\cdot w_i^{(\hat\beta^{\backslash i})} & \leq \frac{1}{\sum_{j}w_j^{(\beta^{\backslash j})}}\sum_{i=0}^{t}(l_i\cdot w_i^{(\hat\beta^{\backslash t})} + K \epsilon_i)  = \frac{1}{\sum_{j}w_j^{(\beta^{\backslash j})}}\sum_{i=0}^{t}l_i\cdot w_i^{(\hat\beta^{\backslash t})} + \frac{1}{\sum_{j}w_j^{(\beta^{\backslash j})}}\sum_{i=0}^{t}K \epsilon_i
% \end{align*}
% }

% {\color{blue}
% Assuming that the \textit{unnormalized} weights are $K$-Lipschitz:
% \begin{align*}
%     \frac{1}{\sum_{j}w_j^{(\beta^{\backslash j})}}\sum_{i=0}^{t}l_i\cdot w_i^{(\hat\beta^{\backslash i})}  \leq \frac{1}{\sum_{j}w_j^{(\beta^{\backslash j})}}\sum_{i=0}^{t}l_i\cdot (w_i^{(\hat\beta^{\backslash t})} + K \epsilon_i)  = & \frac{1}{\sum_{j}w_j^{(\beta^{\backslash j})}}\sum_{i=0}^{t}l_i\cdot w_i^{(\hat\beta^{\backslash t})} + \frac{1}{\sum_{j}w_j^{(\beta^{\backslash j})}}\sum_{i=0}^{t}l_iK \epsilon_i \\
%     % \leq & \sum_{i=0}^{t}l_i\cdot \tilde{w}_i^{(\hat\beta^{\backslash t})} + \sum_{i=0}^{t}BK \epsilon_i
% \end{align*}
% \begin{align*}
%     \frac{1}{\sum_{j}w_j^{(\beta^{\backslash j})}}\sum_{i=0}^{t}l_i\cdot w_i^{(\hat\beta^{\backslash i})}  \leq \frac{1}{\sum_{j}w_j^{(\beta^{\backslash j})}}\sum_{i=0}^{t}l_i\cdot (w_i^{(\hat\beta^{\backslash t})} + K \epsilon_i)  = & \frac{1}{\sum_{j}w_j^{(\beta^{\backslash j})}}\sum_{i=0}^{t}l_i\cdot w_i^{(\hat\beta^{\backslash t})} + \frac{1}{\sum_{j}w_j^{(\beta^{\backslash j})}}\sum_{i=0}^{t}l_iK \epsilon_i \\
%     % \leq & \sum_{i=0}^{t}l_i\cdot \tilde{w}_i^{(\hat\beta^{\backslash t})} + \sum_{i=0}^{t}BK \epsilon_i
% \end{align*}
% }


% {\color{blue}
% Assuming that the \textit{normalized} weights are $K$-Lipschitz:
% \begin{align*}
%     \sum_{i=0}^{t}l_i\cdot \tilde{w}_i^{(\hat\beta^{\backslash i})}  \leq \sum_{i=0}^{t}l_i\cdot (\tilde{w}_i^{(\hat\beta^{\backslash t})} + K \epsilon_i)  = & \sum_{i=0}^{t}l_i\cdot \tilde{w}_i^{(\hat\beta^{\backslash t})} + \sum_{i=0}^{t}l_iK \epsilon_i \\
%     \leq & \sum_{i=0}^{t}l_i\cdot \tilde{w}_i^{(\hat\beta^{\backslash t})} + \sum_{i=0}^{t}BK \epsilon_i
% \end{align*}
% \begin{align*}
%     \sum_{i=0}^{t}l_i\cdot \tilde{w}_i^{(\hat\beta^{\backslash i})}  \leq \sum_{i=0}^{t}l_i\cdot (\tilde{w}_i^{(\hat\beta^{\backslash t})} + K \epsilon_i)  = & \sum_{i=0}^{t}l_i\cdot \tilde{w}_i^{(\hat\beta^{\backslash t})} + \sum_{i=0}^{t}l_iK \epsilon_i \\
%     \leq & \sum_{i=0}^{t}l_i\cdot \tilde{w}_i^{(\hat\beta^{\backslash t})} + \sum_{i=0}^{t}BK \epsilon_i
% \end{align*}
% }

% By Eq. \eqref{eq:cpc_pf_step_3_result}, we know that $\hat{\beta}^{\backslash t}=\hat{\beta}\big(\lbag l_{0:t-1}, B \rbag, \tilde{w}_{0:t-1}^{(\cdot)}, \tilde{w}_{\max}^{(\cdot)}, \alpha\big) \leq \beta^*$, and because this $\hat{\beta}^{\backslash t}$ does not depend on $i$, we can apply Eq. \eqref{eq:cpc_step_2} to bound $\sum_{i=0}^{t}l_i\cdot \tilde{w}_i^{(\hat{\beta}^{\backslash t})}$ above by $\alpha$, which gives 
% \begin{align*}
%      \sum_{i=0}^{t}l_i\cdot\tilde{w}_i^{(\hat{\beta}^{\backslash i})} & \leq \alpha + \sum_{i=0}^{t}\frac{K \epsilon_i}{\sum_{j=0}^tw_j^{(\hat\beta^{\backslash t})}}  \\
%      \iff \mathbb{E}\big[L_{t}(\hat\beta) \mid \lbag Z_{0:t} \rbag =  \lbag z_{0:t} \rbag\big] & \leq \alpha + \sum_{i=0}^{t}\frac{K \epsilon_i}{\sum_{j=0}^tw_j^{(\hat\beta^{\backslash t})}}.
% \end{align*}
% Note that $\frac{1}{\sum_{j=0}^tw_j^{(\hat\beta^{\backslash t})}}$ is a normalization constant that plays the role of $\frac{1}{n+1}$ in the case of exchangeable data.
% Then, marginalizing over the event and applying $\epsilon$-replace-one stability, we have
% \begin{align*}
%     \mathbb{E}_{\pi_{0:t}^{(\hat\beta)}}\big[L_{t} \big] & \leq \alpha + K\cdot \epsilon .
% \end{align*}

% Moreover, if the Lipschitz constant $K$ is known, then we can find a conservative empirical significance level, $\hat\alpha$, that can be implemented to control the risk at the target level $\alpha$. To do so, define
% \begin{align*}
%     \hat\epsilon_i = \max_{b\in \{0, B\}}\Big|\hat{\beta}\big(\lbag l_{0:t-1}, B \rbag, \tilde{w}_{0:t-1}^{(\cdot)}, \tilde{w}_{\max}^{(\cdot)}, \alpha\big)-\hat{\beta}\big(\lbag l_{0:t-1\backslash i}, b, B \rbag, \tilde{w}_{0:t-1}^{(\cdot)}, \tilde{w}_{\max}^{(\cdot)}, \alpha\big)\Big|,
% \end{align*}
% which is a conservative estimate of $\epsilon_i$ using only the observations $l_{0:t-1}$, as note that by construction, $\hat\epsilon_i\geq \epsilon_i$. Then, fit the conservative significance level as 
% \begin{align*}
%     \hat\alpha := \sup\Big(\Big\{\alpha' \leq \alpha \quad : \quad \alpha' + \frac{K }{\sum_{j=0}^tw_j^{(\hat\beta^{\backslash t})}}\sum_{i=0}^{t}\epsilon_i(\alpha')\leq \alpha\Big\}  \Big), 
% \end{align*}
% and by a similar argument as above, if $\hat\alpha$ exists,
% \begin{align*}
%     \mathbb{E}_{\pi_{0:t}^{(\hat\beta(\hat\alpha))}}\big[L_{t} \big] & \leq \alpha.
% \end{align*}


% {\color{blue}

% \begin{align*}
%     l_iw_ir_i \leq l_iw_i e^{K\epsilon_i}
% \end{align*}

% Note that 
% \begin{align*}
% r_i = \frac{w_i^{(\hat\beta^{\backslash i})}}{w_i^{(\hat\beta^{\backslash t})}} = \frac{w_i^{(\hat\beta^{\backslash i})} + w_i^{(\hat\beta^{\backslash t})} - w_i^{(\hat\beta^{\backslash t})}}{w_i^{(\hat\beta^{\backslash t})}} = 1 + \frac{w_i^{(\hat\beta^{\backslash i})} - w_i^{(\hat\beta^{\backslash t})}}{w_i^{(\hat\beta^{\backslash t})}}
% \end{align*}

% So, 
% \begin{align*}
%     l_iw_i^{(\hat\beta^{\backslash t})}r_i  & = l_iw_i^{(\hat\beta^{\backslash t})} + l_iw_i^{(\hat\beta^{\backslash t})}\frac{w_i^{(\hat\beta^{\backslash i})} - w_i^{(\hat\beta^{\backslash t})}}{w_i^{(\hat\beta^{\backslash t})}} \\
%     & = l_iw_i^{(\hat\beta^{\backslash t})} + l_i(w_i^{(\hat\beta^{\backslash i})} - w_i^{(\hat\beta^{\backslash t})}) \\
%     & \leq l_iw_i^{(\hat\beta^{\backslash t})} + K\epsilon_i
% \end{align*}

% \begin{align*}
%     l_iw_i^{(\hat\beta^{\backslash t})} - K\epsilon_i \leq l_iw_i^{(\hat\beta^{\backslash t})}r_i \leq l_iw_i^{(\hat\beta^{\backslash t})} + K\epsilon_i
% \end{align*}


% So, 
% \begin{align*}
%     l_iw_i^{(\hat\beta^{\backslash t})}r_{\sigma'(i)}  & = l_iw_i^{(\hat\beta^{\backslash t})} + l_iw_i^{(\hat\beta^{\backslash t})}\frac{w_{\sigma'(i)}^{(\hat\beta^{\backslash i})} - w_{\sigma'(i)}^{(\hat\beta^{\backslash t})}}{w_{\sigma'(i)}^{(\hat\beta^{\backslash t})}} \\
%     & = l_iw_i^{(\hat\beta^{\backslash t})} + l_iw_i^{(\hat\beta^{\backslash t})}\frac{w_{\sigma'(i)}^{(\hat\beta^{\backslash i})} - w_{\sigma'(i)}^{(\hat\beta^{\backslash t})}}{w_{\sigma'(i)}^{(\hat\beta^{\backslash t})}}
% \end{align*}


% \begin{align*}
%     \mathbb{E}[\epsilon_i \mid \lbag l_{0:t} \rbag] = \epsilon_i\cdot\mathbb{P}(Z_t = z_i  \mid \lbag l_{0:t} \rbag) = \epsilon_i \cdot \tilde{w}_i^{(\hat\beta^{\backslash t})} = \epsilon_i \cdot \frac{w_i^{(\hat\beta^{\backslash t})}}{\sum_j w_j^{(\hat\beta^{\backslash t})}}
% \end{align*}


% }

% Alternative last analysis:

\begin{align}\label{eq:cpc_pf_branch_point}
    \frac{\sum_{j}w_j^{(\hat\beta^{\backslash t})}}{\sum_{j}w_j^{(\hat\beta^{\backslash j})}}\cdot\alpha & +  \frac{1}{\sum_{j}w_j^{(\hat\beta^{\backslash j})}}\sum_{i=0}^{t}(l_i\cdot w_i^{(\hat\beta^{\backslash i})} - l_i\cdot w_i^{(\hat\beta^{\backslash t})}) \nonumber \\
    % = \cdots \\
    = & \frac{\sum_{j}(w_j^{(\hat\beta^{\backslash t})} + w_j^{(\hat\beta^{\backslash j})} - w_j^{(\hat\beta^{\backslash j})})}{\sum_{j}w_j^{(\hat\beta^{\backslash j})}}\cdot\alpha +  \frac{1}{\sum_{j}w_j^{(\hat\beta^{\backslash j})}}\sum_{i=0}^{t}(l_i\cdot w_i^{(\hat\beta^{\backslash i})} - l_i\cdot w_i^{(\hat\beta^{\backslash t})}) \nonumber \\
    % = & \frac{\sum_{j}w_j^{(\hat\beta^{\backslash j})} + \sum_{j}(w_j^{(\hat\beta^{\backslash t})} - w_j^{(\hat\beta^{\backslash j})})}{\sum_{j}w_j^{(\hat\beta^{\backslash j})}}\cdot\alpha +  \frac{1}{\sum_{j}w_j^{(\hat\beta^{\backslash j})}}\sum_{i=0}^{t}(l_i\cdot w_i^{(\hat\beta^{\backslash i})} - l_i\cdot w_i^{(\hat\beta^{\backslash t})}) \\
    % = & \bigg(1 + \frac{\sum_{j}(w_j^{(\hat\beta^{\backslash t})} - w_j^{(\hat\beta^{\backslash j})})}{\sum_{j}w_j^{(\hat\beta^{\backslash j})}}\bigg)\cdot\alpha +  \frac{1}{\sum_{j}w_j^{(\hat\beta^{\backslash j})}}\sum_{i=0}^{t}(l_i\cdot w_i^{(\hat\beta^{\backslash i})} - l_i\cdot w_i^{(\hat\beta^{\backslash t})}) \\ 
    = & \ \alpha + \frac{\sum_{j}(w_j^{(\hat\beta^{\backslash t})} - w_j^{(\hat\beta^{\backslash j})})\cdot\alpha}{\sum_{j}w_j^{(\hat\beta^{\backslash j})}} +  \frac{\sum_{i=0}^{t}(l_i\cdot w_i^{(\hat\beta^{\backslash i})} - l_i\cdot w_i^{(\hat\beta^{\backslash t})})}{\sum_{j}w_j^{(\hat\beta^{\backslash j})}} \nonumber \\
    % = & \ \alpha + \frac{\sum_{i=0}^{t}(\alpha\cdot w_i^{(\hat\beta^{\backslash t})} - \alpha\cdot w_i^{(\hat\beta^{\backslash i})}) + \sum_{i=0}^{t}(l_i\cdot w_i^{(\hat\beta^{\backslash i})} - l_i\cdot w_i^{(\hat\beta^{\backslash t})})}{\sum_{j}w_j^{(\hat\beta^{\backslash j})}} \\
    = & \ \alpha + \frac{\sum_{i=0}^{t}\alpha\cdot w_i^{(\hat\beta^{\backslash t})} - \alpha\cdot w_i^{(\hat\beta^{\backslash i})} + l_i\cdot w_i^{(\hat\beta^{\backslash i})} - l_i\cdot w_i^{(\hat\beta^{\backslash t})}}{\sum_{j}w_j^{(\hat\beta^{\backslash j})}} \nonumber \\
    % = & \ \alpha + \frac{\sum_{i=0}^{t}l_i\cdot (w_i^{(\hat\beta^{\backslash i})} - w_i^{(\hat\beta^{\backslash t})}) + \alpha\cdot (w_i^{(\hat\beta^{\backslash t})} -  w_i^{(\hat\beta^{\backslash i})})}{\sum_{j}w_j^{(\hat\beta^{\backslash j})}} \\
    = & \ \alpha + \frac{\sum_{i=0}^{t}l_i\cdot (w_i^{(\hat\beta^{\backslash i})} - w_i^{(\hat\beta^{\backslash t})}) - \alpha\cdot (w_i^{(\hat\beta^{\backslash i})} -  w_i^{(\hat\beta^{\backslash t})})}{\sum_{j}w_j^{(\hat\beta^{\backslash j})}} \nonumber \\
    = & \ \alpha + \frac{\sum_{i=0}^{t}(l_i - \alpha)\cdot (w_i^{(\hat\beta^{\backslash i})} -  w_i^{(\hat\beta^{\backslash t})})}{\sum_{j}w_j^{(\hat\beta^{\backslash j})}}.
\end{align}



Note that $(l_i - \alpha) \leq B$, so we obtain the inequality

\begin{align}\label{eq:cpc_pf_b_sub}
    \frac{\sum_{j}w_j^{(\hat\beta^{\backslash t})}}{\sum_{j}w_j^{(\hat\beta^{\backslash j})}}\cdot\alpha & +  \frac{1}{\sum_{j}w_j^{(\hat\beta^{\backslash j})}}\sum_{i=0}^{t}(l_i\cdot w_i^{(\hat\beta^{\backslash i})} - l_i\cdot w_i^{(\hat\beta^{\backslash t})}) \quad \leq \quad  \ \alpha + B\frac{\sum_{i=0}^{t}(w_i^{(\hat\beta^{\backslash i})} -  w_i^{(\hat\beta^{\backslash t})})}{\sum_{j}w_j^{(\hat\beta^{\backslash j})}}.
\end{align}

Assuming that $w_i^{(\beta)}$ is $K$-Lipschitz in $\beta$, we can bound the RHS of Eq. \eqref{eq:cpc_pf_b_sub} as

\begin{align}\label{eq:cpc_pf_b_sub_lipschitz}
    \alpha + B\frac{\sum_{i=0}^{t}(w_i^{(\hat\beta^{\backslash i})} -  w_i^{(\hat\beta^{\backslash t})})}{\sum_{j}w_j^{(\hat\beta^{\backslash j})}} \quad \leq \quad  \ \alpha + B\frac{\sum_{i=0}^{t}K\epsilon_i}{\sum_{j}w_j^{(\hat\beta^{\backslash j})}}.
\end{align}

Using transitive property across Inequality \eqref{eq:risk_bound_alpha_sub}, Equation \eqref{eq:cpc_pf_branch_point}, Inequality \eqref{eq:cpc_pf_b_sub}, and Inequality \eqref{eq:cpc_pf_b_sub_lipschitz}, we have

\begin{align*}
    \mathbb{E}_{\pi_{0:t}^{(\hat\beta)}}\Big[L_{t} \mid \lbag Z_{0:t} \rbag = \lbag z_{0:t} \rbag\Big] \leq \alpha + BK\frac{\sum_{i=0}^{t}\epsilon_i}{\sum_{j}w_j^{(\hat\beta^{\backslash j})}},
\end{align*}

and marginalizing over the event $\lbag Z_{0:t} \rbag = \lbag z_{0:t} \rbag$ gives 

\begin{align}\label{cpc_pf_result_1}
    \mathbb{E}_{\pi_{0:t}^{(\hat\beta)}}\big[L_{t} \big] \leq \alpha + BK\epsilon.
\end{align}


Alternatively, if from Eq. \eqref{eq:cpc_pf_branch_point} we instead assumed that $(l_i - \alpha) \cdot w_i^{(\beta)}$ is $K$ Lipschitz in $\beta$, then from Eq. \eqref{eq:cpc_pf_branch_point} we obtain the inequality
\begin{align}\label{eq:cpc_pf_li_alpha_lipschitz}
    \alpha + \frac{\sum_{i=0}^{t}(l_i - \alpha)(w_i^{(\hat\beta^{\backslash i})} -  w_i^{(\hat\beta^{\backslash t})})}{\sum_{j}w_j^{(\hat\beta^{\backslash j})}} \quad \leq \quad  \ \alpha + \frac{\sum_{i=0}^{t}K\epsilon_i}{\sum_{j}w_j^{(\hat\beta^{\backslash j})}},
\end{align}

and using transitive property across Inequality \eqref{eq:risk_bound_alpha_sub}, Equation \eqref{eq:cpc_pf_branch_point}, and Inequality \eqref{eq:cpc_pf_li_alpha_lipschitz}, we have

\begin{align*}
    \mathbb{E}_{\pi_{0:t}^{(\hat\beta)}}\Big[L_{t} \mid \lbag Z_{0:t} \rbag = \lbag z_{0:t} \rbag\Big] \leq \alpha + K\frac{\sum_{i=0}^{t}\epsilon_i}{\sum_{j}w_j^{(\hat\beta^{\backslash j})}},
\end{align*}

and marginalizing over the event $\lbag Z_{0:t} \rbag = \lbag z_{0:t} \rbag$ gives 

\begin{align}\label{cpc_pf_result_2}
    \mathbb{E}_{\pi_{0:t}^{(\hat\beta)}}\big[L_{t} \big] \leq \alpha + K\epsilon.
\end{align}



% Assume that $(l_i - \alpha)\cdot w_i^{(\beta)}$ is $K$-Lipschitz, we have
% \begin{align*}
%     |(l_i - \alpha)\cdot (w_i^{(\hat\beta^{\backslash i})} -  w_i^{(\hat\beta^{\backslash t})})| \leq K\epsilon_i \\
%     \implies (l_i - \alpha)\cdot (w_i^{(\hat\beta^{\backslash i})} -  w_i^{(\hat\beta^{\backslash t})}) \leq K\epsilon_i.
% \end{align*}
% Note that we can roughly interpret this as assuming that the oracle weight we would get by replacing one sample, $w_i^{(\hat\beta^{\backslash i})}$, does not put too much \textit{more} weight than the empirical version, $w_i^{(\hat\beta^{\backslash t})}$, on losses that exceed $\alpha$ (i.e., that replacing one sample doesn't greatly underestimate the weight on ``unsafe'' losses); and, that the oracle weight, $w_i^{(\hat\beta^{\backslash i})}$, does not put far \textit{less} weight than the empirical version, $w_i^{(\hat\beta^{\backslash t})}$, on losses that are under $\alpha$ (i.e., that replacing one sample doesn't greatly overestimate the weight on ``safe'' losses). 

% Then,
% \begin{align*}
%     \frac{\sum_{j}w_j^{(\hat\beta^{\backslash t})}}{\sum_{j}w_j^{(\hat\beta^{\backslash j})}}\cdot\alpha & +  \frac{1}{\sum_{j}w_j^{(\hat\beta^{\backslash j})}}\sum_{i=0}^{t}(l_i\cdot w_i^{(\hat\beta^{\backslash i})} - l_i\cdot w_i^{(\hat\beta^{\backslash t})}) \leq  &  \alpha + \frac{\sum_{i=0}^{t}K\epsilon_i}{\sum_{j}w_j^{(\hat\beta^{\backslash j})}} 
% \end{align*}

% and so

% \begin{align*}
%     \mathbb{E}[L_t] \leq \alpha + K\epsilon. 
% \end{align*}

% Alternatively, assume that $\max(l_i, B-l_i)\cdot w_i^{(\beta)}$ is $K$-Lipschitz in $\beta$, and the above analysis should also follow


% {
% \color{blue}

% \begin{align*}
%     % \frac{\sum_{i=0}^{t}(l_i - \alpha)\cdot (w_i^{(\hat\beta^{\backslash i})} -  w_i^{(\hat\beta^{\backslash t})})}{\sum_{j}w_j^{(\hat\beta^{\backslash j})}} = 
%     & \frac{\sum_{i=0}^{t}(l_i - \alpha)  \cdot (w_i^{(\hat\beta^{\backslash i})} -  w_i^{(\hat\beta^{\backslash t})})}{\sum_{j}w_j^{(\hat\beta^{\backslash j})}} \\
%     & = \frac{\sum_{i=0}^{t}l_i \cdot w_i^{(\hat\beta^{\backslash i})}}{\sum_{j}w_j^{(\hat\beta^{\backslash j})}} - \frac{\sum_{i=0}^{t}l_i \cdot w_i^{(\hat\beta^{\backslash t})}}{\sum_{j}w_j^{(\hat\beta^{\backslash j})}} + \frac{\sum_{i=0}^{t}\alpha \cdot w_i^{(\hat\beta^{\backslash i})}}{\sum_{j}w_j^{(\hat\beta^{\backslash j})}} - \frac{\sum_{i=0}^{t}\alpha \cdot w_i^{(\hat\beta^{\backslash t})}}{\sum_{j}w_j^{(\hat\beta^{\backslash j})}} \\
%     & = \mathbb{E}[L_t \mid \lbag Z_{0:t} \rbag = \lbag z_{0:t} \rbag] - \frac{\sum_{i=0}^{t}l_i \cdot w_i^{(\hat\beta^{\backslash t})}}{\sum_{j}w_j^{(\hat\beta^{\backslash j})}} + \mathbb{E}[\ \alpha \mid \lbag Z_{0:t} \rbag = \lbag z_{0:t} \rbag] + \frac{\sum_{i=0}^{t}\alpha \cdot w_i^{(\hat\beta^{\backslash t})}}{\sum_{j}w_j^{(\hat\beta^{\backslash j})}}
% \end{align*}

% \begin{align*}
%     \frac{\sum_{i=0}^{t}l_i \cdot w_i^{(\hat\beta^{\backslash t})}}{\sum_{j}w_j^{(\hat\beta^{\backslash j})}} 
%     & = \frac{l_i\cdot\sum_{\sigma:\sigma(t)=i}\pi_{0:t}^{(\hat\beta^{\backslash t})}(z_{\sigma(0)}, ..., z_{\sigma(t)})}{\sum_{\sigma}\pi_{0:t}^{(\hat\beta)}(z_{\sigma(0)}, ..., z_{\sigma(t)})} \\
%     & = \frac{l_i\cdot\sum_{\sigma:\sigma(t)=i}\pi_{0:t}^{(\hat\beta^{\backslash t})}(z_{\sigma(0)}, ..., z_{\sigma(t)})}{\sum_{\sigma}\pi_{0:t}^{(\hat\beta)}(z_{\sigma(0)}, ..., z_{\sigma(t)})} \\
%     & = \frac{l_i\cdot\mathbb{P}_{\pi_{0:t}^{(\hat\beta^{\backslash t})}}\big(Z_{t}=z_i, \  \lbag Z_{0:t} \rbag = \lbag z_{0:t} \rbag\big)}{\mathbb{P}_{\pi_{0:t}^{(\hat\beta)}}\big(\lbag Z_{0:t} \rbag = \lbag z_{0:t} \rbag\big)}
% \end{align*}


% \begin{align*}
%     \frac{\sum_{i=0}^{t}l_i  \cdot (w_i^{(\hat\beta^{\backslash i})} -  w_i^{(\hat\beta^{\backslash t})})}{\sum_{j}w_j^{(\hat\beta^{\backslash j})}} &  = \sum_{i=0}^{t}l_i \cdot \frac{w_i^{(\hat\beta^{\backslash i})}}{\sum_{j}w_j^{(\hat\beta^{\backslash j})}} - \sum_{i=0}^{t}l_i \cdot \frac{ w_i^{(\hat\beta^{\backslash t})}}{\sum_{j}w_j^{(\hat\beta^{\backslash j})}} \\
%     & = \sum_{i=0}^{t}l_i\cdot\frac{\mathbb{P}_{\pi_{0:t}^{(\hat\beta^{\backslash i})}}\big(Z_{t}=z_i, \  \lbag Z_{0:t} \rbag = \lbag z_{0:t} \rbag\big)}{\mathbb{P}_{\pi_{0:t}^{(\hat\beta)}}\big(\lbag Z_{0:t} \rbag = \lbag z_{0:t} \rbag\big)} - \sum_{i=0}^{t}l_i\cdot\frac{\mathbb{P}_{\pi_{0:t}^{(\hat\beta^{\backslash t})}}\big(Z_{t}=z_i, \  \lbag Z_{0:t} \rbag = \lbag z_{0:t} \rbag\big)}{\mathbb{P}_{\pi_{0:t}^{(\hat\beta)}}\big(\lbag Z_{0:t} \rbag = \lbag z_{0:t} \rbag\big)}
% \end{align*}

% Note that marginalizing effectively means multiplying by the denominator $\mathbb{P}_{\pi_{0:t}^{(\hat\beta)}}\big(\lbag Z_{0:t} \rbag = \lbag z_{0:t} \rbag\big)$ and summing over the event $\lbag Z_{0:t} \rbag = \lbag z_{0:t} \rbag$.

% }

% \vspace{2cm}

% \textbf{Attempt trying to remove $\alpha$, assuming $l_i\cdot w_i^{(\beta)}$ is Lipschitz:} Partition the indices based on the sign of $l_i -\alpha$ (i.e., based on whether $l_i \geq \alpha$):
% \begin{align*}
%     \mathcal{I}_+ & := \big\{i\in \{0, ..., t\} : l_i \geq \alpha\big\} \\
%     \mathcal{I}_- & := \big\{i\in \{0, ..., t\} : l_i < \alpha \big\},
% \end{align*}

% then, 
% \begin{align*}
%     \text{RHS} & = \ \alpha + \frac{\sum_{i\in\mathcal{I}_+}(l_i -\alpha)(w_i^{(\hat\beta^{\backslash i})} - w_i^{(\hat\beta^{\backslash t})}) + \sum_{i\in\mathcal{I}_-}(l_i -\alpha)(w_i^{(\hat\beta^{\backslash i})} - w_i^{(\hat\beta^{\backslash t})})}{\sum_{j}w_j^{(\hat\beta^{\backslash j})}} \\
%     & = \ \alpha + \frac{\sum_{i\in\mathcal{I}_+}|l_i -\alpha|(w_i^{(\hat\beta^{\backslash i})} - w_i^{(\hat\beta^{\backslash t})}) - \sum_{i\in\mathcal{I}_-}|l_i -\alpha|(w_i^{(\hat\beta^{\backslash i})} - w_i^{(\hat\beta^{\backslash t})})}{\sum_{j}w_j^{(\hat\beta^{\backslash j})}} \\
%     & \leq \ \alpha + \frac{\sum_{i\in\mathcal{I}_+}|l_i -\alpha|\cdot |w_i^{(\hat\beta^{\backslash i})} - w_i^{(\hat\beta^{\backslash t})}| + \sum_{i\in\mathcal{I}_-}|l_i -\alpha|\cdot |w_i^{(\hat\beta^{\backslash i})} - w_i^{(\hat\beta^{\backslash t})}|}{\sum_{j}w_j^{(\hat\beta^{\backslash j})}} \\
%     % & \leq \ \alpha + \frac{\sum_{i\in\mathcal{I}_+}l_i \cdot |w_i^{(\hat\beta^{\backslash i})} - w_i^{(\hat\beta^{\backslash t})}| + \sum_{i\in\mathcal{I}_-}l_i\cdot |w_i^{(\hat\beta^{\backslash i})} - w_i^{(\hat\beta^{\backslash t})}|}{\sum_{j}w_j^{(\hat\beta^{\backslash j})}} 
% \end{align*}



\end{proof}


\subsection{Counterexample to gCRC risk control on arbitrary bounded losses}
\label{app:counterexample}


\begin{proposition}
    
    The following assumptions are not sufficient for $\hat\lambda_+$ as defined in Eq. \eqref{eq:methods_lambda_empirical_def} to achieve the risk control guarantee at level $\alpha$:
    \begin{itemize}
        \item Right-continuity of $L_i(\lambda)$ in $\lambda$,
        \item Existence of a safe hyperparameter: For $\lambda_{\max}:=\sup\Lambda\in \Lambda$, that $L_i(\lambda_{\max})\leq \alpha$,
        \item Boundedness of the loss functions: $\sup_{\lambda}L_i(\lambda)\leq B < \infty \quad \text{almost surely}$.
    \end{itemize}
    
    
\end{proposition}


\begin{proof}
This counterexample assumes exchangeable data.
Let $\Lambda = \{\lambda_1, \lambda_2, \lambda_3, \lambda_{\max}\}$ where $\lambda_1 < \lambda_2 < \lambda_3 < \lambda_{\max}$.
Fix $\alpha \in (0, 1]$.
Consider the following $n + 1 = 3$ loss functions, where $B = \frac{3}{2}\alpha$:
\begin{align}
\label{eq:bag-loss}
l_1(\lambda) =
    \begin{cases}
    B, & \lambda = \lambda_1 \\
    B/2, & \lambda = \lambda_2 \\
    0, & \lambda = \lambda_3 \\
    0, & \lambda = \lambda_{\max}
  \end{cases}, \qquad 
    l_2(\lambda) =
    \begin{cases}
    B/2, & \lambda = \lambda_1 \\
    B, & \lambda = \lambda_2 \\
    0, & \lambda = \lambda_3 \\
    0, & \lambda = \lambda_{\max}
  \end{cases}, \qquad 
    l_3(\lambda) =
    \begin{cases}
    B/2, & \lambda = \lambda_1 \\
    0, & \lambda = \lambda_2 \\
    B, & \lambda = \lambda_3 \\
    0, & \lambda = \lambda_{\max}
  \end{cases}.
\end{align}
Note that $\max_{\lambda \in \Lambda} l_i(\lambda) = B$ and $l_i(\lambda_{\max}) = 0 < \alpha$ for all $i \in [3]$, satisfying our assumptions.

Consider the bag-conditional empirical risk,
\begin{align}
\label{eq:tower-again}
    \mathbb{E}[L_3(\hat{\lambda}) \mid \lbag L_{1:3}\rbag = \lbag l_{1:3}\rbag] & = \mathbb{E}\left[ \; \mathbb{E}[ L_3(\hat{\lambda}) \mid \lbag L_{1:3}\rbag = \lbag l_{1:3}\rbag, L_3 = l_i] \; \middle| \; \lbag L_{1:3}\rbag = \lbag l_{1:3}\rbag \right],
\end{align}
where the inner expectation on the righthand side conditions on the value of the bag of loss functions, $\lbag L_{1:3} \rbag = \lbag l_{1:3} \rbag$, and additionally on the event that the test loss takes the value of the $i$-th loss function in the bag, $L_{3} = l_i$.
Note that the test loss at the empirical threshold, $L_{3}(\hat{\lambda})$, is a deterministic function of the values of $\lbag L_{1:3} \rbag$ and $L_{3}$.
Therefore, this inner expectation is simply
\begin{align}
    \mathbb{E}\left[L_{3}(\hat{\lambda}) \; \middle| \; \lbag L_{1:3}\rbag = \lbag l_{1:3}\rbag, L_{3} = l_i\right] & = l_i(\hat{\lambda}(\lbag l_{1:3\setminus i}, B \rbag, \alpha)).
\end{align}
Substituting this expression for the inner expectation in Eq.~\eqref{eq:tower-again}, we can write the bag-conditional empirical risk as
\begin{align}
\label{eq:bag-conditional-empirical-risk}
    \mathbb{E}[L_3(\hat{\lambda}) \mid \lbag L_{1:3}\rbag = \lbag l_{1:3}\rbag] = \mathbb{E}[ \; l_i(\hat{\lambda}(\lbag l_{1:3\setminus i}, B \rbag, \alpha)) \mid \lbag L_{1:3}\rbag = \lbag l_{1:3}\rbag] = \frac{1}{3}\sum_{i =1}^{3} l_i(\hat{\lambda}(\lbag l_{1:3\setminus i}, B \rbag, \alpha)).
\end{align}
We now calculate this quantity for the bag of loss functions given in Eqs.~\eqref{eq:bag-loss}.
Each case or ``world'' below computes one summand in the mean in Eq.~\eqref{eq:bag-conditional-empirical-risk}.
To simplify forthcoming calculations, note that
\begin{align*}
    \hat{\lambda}(\lbag l_{1:3\setminus i}, B \rbag, \alpha) = \min \{\lambda' \in \Lambda : \forall \lambda \geq \lambda', \sum_{j \neq i}l_j(\lambda) \leq 3\alpha - B = 2B - B = B\}.
\end{align*}
\begin{itemize}
    \item \textbf{Case 1:} $L_3 = l_1$. We have
    \begin{align}
        \hat{\lambda}(\lbag l_2, l_3, B \rbag, \alpha) & = \min \{\lambda' \in \Lambda : \forall \lambda \geq \lambda', l_2(\lambda) + l_3(\lambda) \leq B\} \\
        & = \lambda_1,
    \end{align}
    so $L_3(\hat{\lambda}) = l_1(\hat{\lambda}(\lbag l_2, l_3, B \rbag, \alpha)) = l_1(\lambda_1) = B$.
    \item \textbf{Case 2:} $L_3 = l_2$. We have
    \begin{align}
        \hat{\lambda}(\lbag l_1, l_3, B \rbag, \alpha) & = \min \{\lambda' \in \Lambda : \forall \lambda \geq \lambda', l_1(\lambda) + l_3(\lambda) \leq B\} \\
        & = \lambda_2,
    \end{align}
    so $L_3(\hat{\lambda}) = l_2(\hat{\lambda}(\lbag l_1, l_3, B \rbag, \alpha)) = l_2(\lambda_2) = B$.
    \item \textbf{Case 3:} $L_3 = l_3$. We have
    \begin{align}
        \hat{\lambda}(\lbag l_1, l_2, B \rbag, \alpha) & = \min \{\lambda' \in \Lambda : \forall \lambda \geq \lambda', l_1(\lambda) + l_2(\lambda) \leq B\} \\
        & = \lambda_3,
    \end{align}
    so $L_3(\hat{\lambda}) = l_3(\hat{\lambda}(\lbag l_1, l_2, B \rbag, \alpha)) = l_3(\lambda_3) = B$.
\end{itemize}
Plugging these values into Eq.~\eqref{eq:bag-conditional-empirical-risk} gives
\begin{align*}
    \mathbb{E}[L_3(\hat{\lambda}) \mid \lbag L_{1:3}\rbag = \lbag l_{1:3}\rbag] & = \frac{1}{3}\sum_{i =1}^{3} l_i(\hat{\lambda}(\lbag l_{1:3\setminus i}, B \rbag, \alpha)) \\
    & = \frac{1}{3}(B + B + B) \\
    & = B = \frac{3}{2}\alpha > \alpha.
\end{align*}

\end{proof}

\subsection{Conservative adjustment to gCRC for level $\alpha$ validity if $K$ is known}
\label{app:conservative_adjustment_K_known}

Assume the gCRC setting of Theorem \ref{thm:gcrc_nonmonotonic} in Section \ref{subsec:finite_sample_guarantees_crc_non_monotone} (although an analogous result would also hold for conformal policy control). In practice, the deviation of the empirical gCRC algorithm due to replacing one sample can be bounded above by
\begin{align*}
    \hat\epsilon_i := \max_{b\in \{0, B\}}\Big|\hat\lambda_+(\lbag l_{1:n}, B\rbag, \alpha)-\hat\lambda_+(\lbag l_{1:n\backslash i}, b, B\rbag, \alpha)\Big|,
\end{align*}
and if the Lipschitz constant for the loss functions is known, a conservative significance level can be obtained as
\begin{align*}
    \hat\alpha := \sup\Big(\Big\{\alpha' \leq \alpha \quad : \quad \alpha' + \frac{K}{n+1}\sum_{i=1}^{n+1}\hat\epsilon_i(\alpha')\leq \alpha\Big\}  \Big). 
\end{align*}
Then, the procedure run at $\hat\alpha$ attains level $\alpha$ validity.
\begin{proposition}
    In the gCRC setting of Theorem \ref{thm:gcrc_nonmonotonic},
    \begin{align*}
        \mathbb{E}\big[L_{n+1}\big(\hat{\lambda}_{+}(\hat\alpha)\big)\big] \leq \alpha.
    \end{align*}
\end{proposition}


% \subsection{Clarifying $\epsilon$-stability in policy control setting}

% Recall the definition of $\epsilon$-replace-one stability from the main text, reproduced here.

% \begin{definition}[Restatement of Definition \ref{eq:def_replace_one_stability}]
%     We say that a function  $h$ is $\epsilon$-replace-one stable if for all $i, j\in [n+1]$, 
% \begin{align*}
%     \mathbb{E}\Big[\big|h(\lbag L_{1:n+1\backslash i})-h(\lbag l_{1:n+1\backslash j})\big|\Big] \leq \epsilon.
% \end{align*}
% \end{definition}

% It is worth further clarifying the meaning of this stability condition for $\hat{\beta}:=\hat{\beta}(\lbag L_{1:t-1}, B \rbag, \tilde{w}_{0:t-1}^{(\cdot)}, \tilde{w}_{\max}^{(\cdot)},  \alpha)$, which itself depends on the policies or distributions that the expectation is taken over, via the conformal weights $\tilde{w}_{0:t-1}^{(\cdot)}$ and $\tilde{w}_{\max}^{(\cdot)}$. 

% In the expectation in the definition, set $h = \hat\beta$,  $i=t$,  and $j=t$, and then expand via law of total expectation over $\lbag Z_{0:t}\rbag = \lbag z_{0:t}\rbag$, \footnote{Also note that for policy control we shifted indices from $\{1, ..., n+1\}$ to $\{0, ..., t\}$.}
% \begin{align*}
%     \mathbb{E}\Big[&\big|\hat{\beta}(\lbag L_{0:t-1}, B \rbag, \tilde{w}_{0:t-1}^{(\cdot)}, \tilde{w}_{\max}^{(\cdot)},  \alpha)-\hat{\beta}(\lbag l_{0:t-1}, B \rbag, \tilde{w}_{0:t-1}^{(\cdot)}, \tilde{w}_{\max}^{(\cdot)},  \alpha)\big|\Big]  \\
%     & = \mathbb{E}\Big[\mathbb{E}\Big[\big|\hat{\beta}(\lbag L_{0:t-1}, B \rbag, \tilde{w}_{0:t-1}^{(\cdot)}, \tilde{w}_{\max}^{(\cdot)},  \alpha)-\hat{\beta}(\lbag l_{0:t-1}, B \rbag, \tilde{w}_{0:t-1}^{(\cdot)}, \tilde{w}_{\max}^{(\cdot)},  \alpha)\big| \quad \mid \quad \lbag Z_{0:t}\rbag = \lbag z_{0:t}\rbag\Big]\Big] \\
%     & = \mathbb{E}\Big[\mathbb{E}\Big[\mathbb{E}\Big[\big|\hat{\beta}(\lbag L_{0:t-1}, B \rbag, \tilde{w}_{0:t-1}^{(\cdot)}, \tilde{w}_{\max}^{(\cdot)},  \alpha)-\hat{\beta}(\lbag l_{0:t-1}, B \rbag, \tilde{w}_{0:t-1}^{(\cdot)}, \tilde{w}_{\max}^{(\cdot)},  \alpha)\big| \ \mid \ \lbag Z_{0:t}\rbag = \lbag z_{0:t}\rbag, Z_t = z_i\Big]\ \mid \ \lbag Z_{0:t}\rbag = \lbag z_{0:t}\rbag\Big]\Big]  \\
%     & = \mathbb{E}\Big[\mathbb{E}\Big[\mathbb{E}\Big[\big|\hat{\beta}(\lbag l_{0:t\backslash i}, B \rbag, \tilde{w}_{0:t-1}^{(\cdot)}, \tilde{w}_{\max}^{(\cdot)},  \alpha)-\hat{\beta}(\lbag l_{0:t-1}, B \rbag, \tilde{w}_{0:t-1}^{(\cdot)}, \tilde{w}_{\max}^{(\cdot)},  \alpha)\big| \ \mid \ \lbag Z_{0:t}\rbag = \lbag z_{0:t}\rbag, Z_t = z_i\Big]\ \mid \ \lbag Z_{0:t}\rbag = \lbag z_{0:t}\rbag\Big]\Big] 
% \end{align*}

% Noting that $|\hat{\beta}(\lbag l_{0:t\backslash i}, B \rbag, \tilde{w}_{0:t-1}^{(\cdot)}, \tilde{w}_{\max}^{(\cdot)},  \alpha)-\hat{\beta}(\lbag l_{0:t-1}, B \rbag, \tilde{w}_{0:t-1}^{(\cdot)}, \tilde{w}_{\max}^{(\cdot)},  \alpha)\big|$ is constant conditional on $\lbag Z_{0:t}\rbag = \lbag z_{0:t}\rbag, Z_t = z_i$, 
% \begin{align*}
%     \mathbb{E}\Big[&\big|\hat{\beta}(\lbag L_{0:t-1}, B \rbag, \tilde{w}_{0:t-1}^{(\cdot)}, \tilde{w}_{\max}^{(\cdot)},  \alpha)-\hat{\beta}(\lbag l_{0:t-1}, B \rbag, \tilde{w}_{0:t-1}^{(\cdot)}, \tilde{w}_{\max}^{(\cdot)},  \alpha)\big|\Big]  \\
%     % & = \mathbb{E}\Big[\mathbb{E}\Big[\big|\hat{\beta}(\lbag L_{0:t-1}, B \rbag, \tilde{w}_{0:t-1}^{(\cdot)}, \tilde{w}_{\max}^{(\cdot)},  \alpha)-\hat{\beta}(\lbag l_{0:t-1}, B \rbag, \tilde{w}_{0:t-1}^{(\cdot)}, \tilde{w}_{\max}^{(\cdot)},  \alpha)\big| \quad \mid \quad \lbag Z_{0:t}\rbag = \lbag z_{0:t}\rbag\Big]\Big] \\
%     % & = \mathbb{E}\Big[\mathbb{E}\Big[\mathbb{E}\Big[\big|\hat{\beta}(\lbag L_{0:t-1}, B \rbag, \tilde{w}_{0:t-1}^{(\cdot)}, \tilde{w}_{\max}^{(\cdot)},  \alpha)-\hat{\beta}(\lbag l_{0:t-1}, B \rbag, \tilde{w}_{0:t-1}^{(\cdot)}, \tilde{w}_{\max}^{(\cdot)},  \alpha)\big| \ \mid \ \lbag Z_{0:t}\rbag = \lbag z_{0:t}\rbag, Z_t = z_i\Big]\ \mid \ \lbag Z_{0:t}\rbag = \lbag z_{0:t}\rbag\Big]\Big]  \\
%     & = \mathbb{E}\Big[\mathbb{E}\Big[|\hat{\beta}(\lbag l_{0:t\backslash i}, B \rbag, \tilde{w}_{0:t-1}^{(\cdot)}, \tilde{w}_{\max}^{(\cdot)},  \alpha)-\hat{\beta}(\lbag l_{0:t-1}, B \rbag, \tilde{w}_{0:t-1}^{(\cdot)}, \tilde{w}_{\max}^{(\cdot)},  \alpha)\big| \ \mid \ \lbag Z_{0:t}\rbag = \lbag z_{0:t}\rbag\Big]\Big]  \\
%     & = \mathbb{E}\Big[\sum_{i=0}^t|\hat{\beta}(\lbag l_{0:t\backslash i}, B \rbag, \tilde{w}_{0:t-1}^{(\cdot)}, \tilde{w}_{\max}^{(\cdot)},  \alpha)-\hat{\beta}(\lbag l_{0:t-1}, B \rbag, \tilde{w}_{0:t-1}^{(\cdot)}, \tilde{w}_{\max}^{(\cdot)},  \alpha)\big| \ \cdot \mathbb{P}\big(Z_t = z_i\mid \ \lbag Z_{0:t}\rbag = \lbag z_{0:t}\rbag\big)\Big] \\
%     & = \int_{\lbag Z_{0:t}\rbag}\sum_{i=0}^t|\hat{\beta}(\lbag l_{0:t\backslash i}, B \rbag, \tilde{w}_{0:t-1}^{(\cdot)}, \tilde{w}_{\max}^{(\cdot)},  \alpha)-\hat{\beta}(\lbag l_{0:t-1}, B \rbag, \tilde{w}_{0:t-1}^{(\cdot)}, \tilde{w}_{\max}^{(\cdot)},  \alpha)\big| \ \\
%     & \qquad \qquad \qquad \cdot \mathbb{P}\big(Z_t = z_i\mid \ \lbag Z_{0:t}\rbag = \lbag z_{0:t}\rbag\big)\cdot \mathbb{P}\big(\lbag Z_{0:t}\rbag = \lbag z_{0:t}\rbag\big)\\
%     & = \int_{\lbag Z_{0:t}\rbag}\sum_{i=0}^t|\hat{\beta}(\lbag l_{0:t\backslash i}, B \rbag, \tilde{w}_{0:t-1}^{(\cdot)}, \tilde{w}_{\max}^{(\cdot)},  \alpha)-\hat{\beta}(\lbag l_{0:t-1}, B \rbag, \tilde{w}_{0:t-1}^{(\cdot)}, \tilde{w}_{\max}^{(\cdot)},  \alpha)\big| \  \cdot \mathbb{P}\big(Z_t = z_i, \lbag Z_{0:t}\rbag = \lbag z_{0:t}\rbag\big)
% \end{align*}
% Then, for a given fixed $\beta\in (0, \infty]$, let us write out the expectation in the definition more explicitly with respect to our policies, $\pi_{0:t}^{(\beta)} = (\pi_0, \pi_1,  ..., \pi_t^{(\beta)})$, and expand it via law of total expectation, setting $i=t$:

% \begin{align*}
%     \mathbb{E}_{\pi_{0:t}^{(\beta)}}\Big[\big|h(\lbag L_{1:n+1\backslash i})-h(\lbag L_{1:n+1\backslash j})\big|\Big] = 
% \end{align*}





\clearpage
